\documentclass[11pt,oneside]{article}
\usepackage[T1]{fontenc}
\usepackage[utf8]{inputenc}
\usepackage{amsmath,amsthm,mathtools}
\usepackage{newtxtext,newtxmath}
\usepackage[a4paper,margin=28mm,headheight=15pt]{geometry}
\usepackage{microtype,booktabs,longtable,array,enumitem}
\usepackage[dvipsnames]{xcolor}
\usepackage{fancyhdr}
\usepackage{xurl}
\usepackage{hyperref}
\usepackage{bookmark}
\usepackage{listings}
\definecolor{Ink}{HTML}{203B50}
\definecolor{Link}{HTML}{165D83}
\hypersetup{colorlinks=true,linkcolor=Link,citecolor=Link,urlcolor=Link,
 pdftitle={A Single Dilation Recovers Hahn Support},
 pdfauthor={Research article prepared with ChatGPT},
 pdfsubject={Surreal and surcomplex difference fields, support, cohomology, and definability}}
\urlstyle{same}
\setlength{\emergencystretch}{3em}
\setlength{\parskip}{3pt}
\setlength{\parindent}{1.2em}
\pagestyle{fancy}
\fancyhf{}
\fancyhead[L]{\small\color{Ink}A single dilation recovers Hahn support}
\fancyhead[R]{\small Single-source research report}
\fancyfoot[C]{\thepage}
\renewcommand{\headrulewidth}{0.3pt}
\numberwithin{equation}{section}
\newtheorem{theorem}{Theorem}[section]
\newtheorem{lemma}[theorem]{Lemma}
\newtheorem{proposition}[theorem]{Proposition}
\newtheorem{corollary}[theorem]{Corollary}
\theoremstyle{definition}
\newtheorem{definition}[theorem]{Definition}
\newtheorem{example}[theorem]{Example}
\newtheorem{question}[theorem]{Question}
\theoremstyle{remark}
\newtheorem{remark}[theorem]{Remark}
\newcommand{\No}{\mathbf{No}}
\newcommand{\SC}{\No(i)}
\newcommand{\R}{\mathbb R}
\newcommand{\C}{\mathbb C}
\newcommand{\Q}{\mathbb Q}
\newcommand{\Z}{\mathbb Z}
\newcommand{\N}{\mathbb N}
\newcommand{\kk}{\mathbf k}
\newcommand{\K}{\mathcal K}
\newcommand{\OO}{\mathcal O}
\newcommand{\mm}{\mathfrak m}
\newcommand{\M}{\mathcal M}
\newcommand{\id}{\mathrm{id}}
\newcommand{\supp}{\operatorname{supp}}
\newcommand{\Aut}{\operatorname{Aut}}
\newcommand{\Fix}{\operatorname{Fix}}
\newcommand{\im}{\operatorname{im}}
\newcommand{\lc}{\operatorname{lc}}
\newcommand{\Res}{\operatorname{Res}}
\newcommand{\coef}{\operatorname{coef}}
\newcommand{\dd}{\mathsf d}
\newcommand{\hh}{\mathsf h}
\newcommand{\eps}{\varepsilon}
\DeclareMathOperator{\Cen}{Cen}
\DeclareMathOperator{\Nor}{Nor}
\DeclareMathOperator{\Supp}{Supp}
\newcommand{\repo}[1]{\nolinkurl{#1}}% a path in the collection, relative to docs/
\newcommand{\lbl}[1]{\nolinkurl{#1}}% a label of another report, cited by name
\newcommand{\restr}[2]{#1\!\upharpoonright_{#2}}
\newcommand{\pin}{\texttt{465a54b479a1ee842cbf7db1689a7d2f6bfe25e1}}
\lstset{basicstyle=\ttfamily\footnotesize,breaklines=true,columns=fullflexible,
 frame=single,rulecolor=\color{black!20},showstringspaces=false}

\begin{document}
\hypersetup{pageanchor=false}
\begin{titlepage}
\vspace*{4mm}
{\color{Ink}\large SINGLE-SOURCE RESEARCH REPORT\par}
\vspace{2mm}
{\color{Ink}\small Surreal and surcomplex difference fields\par}
\vspace{9mm}
{\color{Ink}\Huge\bfseries A Single Dilation\\[4pt] Recovers Hahn Support\par}
\vspace{8mm}
{\Large Exact difference cohomology, definable normal forms,\\
 and centralizers of monomial automorphisms\par}
\vspace{9mm}
{\large Research article prepared with ChatGPT\\
 for Vladimir Reshetnikov\par}
\vspace{3mm}
{22 September 2026\par}
\vspace{1mm}
{\small Built from one manuscript of batch~22 of the collection,
pinned to repository commit \texttt{465a54b}\par}
\vspace{8mm}
\begin{abstract}
Let $K=\kk((t^\Gamma))$ be a full Hahn field of characteristic zero.
We construct a strong, valuation-nondecreasing partial inverse of
$\sigma-\lambda$ for every coefficient-fixing monomial automorphism
$\sigma(t^g)=\chi(g)t^{\theta g}$. The only obstructions occur at exponents
fixed by $\theta$ on which $\chi$ equals $\lambda$. For finitely many commuting
such automorphisms, these inverses give an explicit deformation retraction
of the entire Koszul complex onto its joint-resonant coefficients. This
computes all additive cohomology groups and the principal-unit part of
multiplicative cohomology, and gives exact criteria for simultaneous
multiplicative equations.

The main structural application is that the single automorphism
$S_2(\sum a_gt^g)=\sum a_gt^{2g}$ defines the coefficient field, constant
coefficient, canonical monomials, every coefficient, the valuation ring,
and truncation in the language of fields with one automorphism. Consequently
its centralizer among \emph{all} field automorphisms consists precisely of
untwisted coefficient--exponent lifts; neither valuation preservation nor
strong additivity needs to be assumed. A uniform interpretation of full
monadic second-order arithmetic yields undecidability, and a coefficient
formula directly witnesses the tree property of the second kind. All
nonidentity rational dilations give the same definable normal-form data
when $\Gamma$ is divisible, although distinct rational dilations are not
conjugate. The elementwise results apply to $\No$ and $\No(i)$ through Conway
normal forms, with proper-class qualifications stated explicitly.
\end{abstract}
\vfill
{\small\color{Ink}\textbf{Status.} Complete proofs are supplied for the stated
results. Novelty is a qualified research claim, not certified priority.
The accompanying finite checks are not a formal verification. No solution
to the published exponential-automorphism questions is asserted.
An AI-assisted research draft, not refereed; no theorem of this report is
Lean-verified. The exponent dilation is written $S_q$, as in the
collection's surcomplex-automorphisms report; the source wrote $D_q$
(Section~\ref{dsup:sec:conventions}). Relations to the collection as it now
stands are in Section~\ref{dsup:sec:position} and
Appendix~\ref{dsup:app:report}. Every label carries the prefix
\texttt{dsup:}.}
\end{titlepage}
\hypersetup{pageanchor=true}

\begingroup
\small
\setlength{\parskip}{0pt}
\tableofcontents
\endgroup
\newpage

\section{Introduction and research position}\label{dsup:sec:intro}

A normal form is often treated as external information attached to a
non-Archimedean field. The underlying field need not visibly identify its
coefficient field, its chosen monomials, or the operation of deleting a
part of a support. This article shows that one particularly simple field
automorphism can recover all of this information.

The automorphism is exponent doubling. It does not square a general series:
\[
 S_2\Bigl(\sum_g a_gt^g\Bigr)=\sum_g a_gt^{2g},
 \qquad
 \Bigl(\sum_g a_gt^g\Bigr)^2=\sum_{g,h}a_ga_ht^{g+h}.
\]
The contrast between these operations is useful. The equation $S_2x=x^2$
recognizes precisely the canonical monomials. A second equation,
$S_2y-y=x-c$, recognizes the constant coefficient $c$ of $x$. Together these
two equations recover arbitrary coefficients. Even the direction of the
valuation is recovered from the constant coefficient of $(1-m)^{-1}$ for a
monomial $m$. This works over $\C$, without ordering the field or naming
complex conjugation.

The difference equation used here admits a considerably more general
solution. We develop that solution first, rather than relying on a special
calculation for doubling. Its extension to commuting automorphisms gives
an explicit cohomology calculation and separates multiplicative
obstructions into a monomial part and a resonant logarithmic part.

\subsection{Principal results}

Theorems~\ref{dsup:thm:resolvent}, \ref{dsup:thm:koszul}, \ref{dsup:thm:definability}, and
\ref{dsup:thm:centralizer} are the main structural statements. Their content is
as follows.

\begin{enumerate}[leftmargin=*,itemsep=4pt]
\item For $\sigma(t^g)=\chi(g)t^{\theta g}$ and $\lambda\ne0$, restriction to
$R_\lambda=\{g:\theta g=g,\ \chi(g)=\lambda\}$ is exactly both the kernel
projection and the cokernel obstruction for $\sigma-\lambda$. A strong
normalized inverse is given by explicit sums along the increasing
halves of exponent orbits.
\item For a commuting $d$-tuple, let $F$ be its common fixed field. There is
an explicit strong contraction of the nonresonant part of the Koszul
complex, and
\[
 H^r(\Z^d,K_{\mathrm{add}})\simeq F^{\binom dr}.
\]
The analogous statement holds for the additive infinitesimal ideal; the
Hahn logarithm turns it into a statement about principal units.
\item In $(K,S_2)$, for nonzero $2$-divisible $\Gamma$, the entire canonical
coefficient and truncation structure is parameter-free definable. In
particular,
\[
 \Cen_{\Aut(K)}(S_2)
 \simeq \Aut(\kk)\times\Aut_{\mathrm{ord}}(\Gamma).
\]
Every member of this centralizer is automatically strong.
\item The same expansion uniformly interprets
$(\N,\mathcal P(\N);+,\in)$ and has an undecidable complete theory. A
definable coefficient relation has $\mathrm{TP}_2$. Over divisible
$\Gamma$, all nonidentity positive rational dilations are interdefinable,
but they are pairwise nonconjugate as field automorphisms.
\end{enumerate}

Here ``strong'' refers to preservation of Hahn summation, not convergence
of finite partial sums. Every assertion about ordinary group cohomology
is initially made for a set-sized Hahn field. Section~\ref{dsup:sec:surreal}
then states exactly what transfers to the proper-class surreal fields.

\subsection{What is inherited and what is proposed as new}

The construction of monomial automorphisms is established Hahn-field
technology, not a contribution of this article. Kaplan--Krapp--Serra
\cite{KKS} study surreal automorphism groups using normal forms. The
repository's surcomplex-automorphism report \cite{RepoAut} also develops
coefficient, exponent, and character lifts. We use this construction but
do not use its more ambitious classification assertions as premises.

The specialization $g\mapsto qg$ lies in Mahler difference algebra.
Faverjon--Roques \cite{FR} treat algorithmic questions for general linear
Mahler equations with Hahn-series solutions. Our first-order,
constant-scalar resolvent does not replace their higher-order variable-
coefficient theory. Its role here is an arbitrary-rank support theorem
and an input to a simultaneous cohomology calculation.

\emph{Related (batch 34).} Linear Mahler equations in a function variable,
with arbitrary-rank Hahn coefficients, now appear in the holonomic-rigidity
report (\repo{surcomplex/holonomic-rigidity-for-entire-hahn-functions}). Its
\lbl{hol:pc:cor:linear} shows that every strongly entire solution of a mixed
differential--Mahler equation $\sum_jL_j\bigl(f(P_j(z))\bigr)=g(z)$ is a
polynomial when the top pullback $P_m$ has degree at least two and larger
than the degrees of the others; with $P_j=z^{b^j}$ this covers linear
Mahler equations. There the substitution acts on the variable of an entire
series, not on the exponents of a scalar, and neither result is used for the
other.

Camacho \cite{Camacho} shows that Hahn fields supplied with truncation or
related normal-form predicates interpret monadic second-order arithmetic.
The arithmetic consequence of knowing support is therefore prior work.
The new proposed step is recovering that support structure from a
\emph{single named dilation}, with no named valuation or monomial group.
We give a concrete Boolean-coefficient coding to make this consequence
transparent. We do not answer Camacho's monomial-definability question in
the smaller language without the extra automorphism.

Multiplicative valued difference fields are also an established subject;
see Pal \cite{Pal}. Our coefficient action is the identity. We make no
claim that these particular structures satisfy the residue difference
closedness hypotheses used in relative quantifier-elimination results.
The obstruction to $\sigma y-y=1$ is already visible in the residue field.

The focused literature and repository review did not locate the complete
orbit-resolvent/Koszul statement, the single-dilation recovery theorem,
or the resulting unrestricted centralizer statement in the form proved
here. These are \emph{proposed additions}, not an assertion of exhaustive
bibliographic priority. No named long-standing open problem is claimed
resolved. In particular, nothing below determines the existence or image
of nontrivial \emph{exponential} automorphisms of $\No$.

\subsection{Repository snapshot and evidence boundary}

The comparison is with the read-only snapshot
\begin{center}\small\pin.\end{center}
The repository catalogue \cite{RepoDocs} lists a substantial collection of
AI-assisted reports; its formalization ledger distinguishes paper proofs
from checked declarations. The root README \cite{Repo} documents strong
Hahn summation, Neumann support control, coefficient maps, and
infinitesimal exponential/logarithm infrastructure. Those are natural
interfaces for formalizing this article.

The catalogue and the surcomplex-automorphism report's guide were inspected,
not every source line in the repository. The repository was not independently
built. No theorem in this article is being represented as already checked
by Lean. Appendix~\ref{dsup:app:verification} records exactly what the delivered
finite verification program does and does not check.

This subsection is the source's, as of its pin. The repository statements
it makes were checked again against the collection as it now stands; the
result, with the one addition that matters here, is in
Appendix~\ref{dsup:app:pinned}.

\subsection{Notation and conventions in this report}\label{dsup:sec:conventions}

This report prints the source manuscript with its notation adapted to the
collection's notation guide \repo{docs/NOTATION.md}. One map was renamed and
one formula symbol was renamed; no normalization changed.

\begin{enumerate}[leftmargin=*,itemsep=3pt]
\item \emph{The exponent dilation is $S_q$.} For $q\in\Q_{>0}$ and divisible
$\Gamma$ ($2$-divisibility suffices for $q=2$),
\[
 S_q\Bigl(\sum_g a_gt^g\Bigr)=\sum_g a_gt^{qg}.
\]
The source wrote $D_q$, and in Section~\ref{dsup:sec:definability} and
Appendix~\ref{dsup:app:formulas} abbreviated $D_2$ to $D$. The notation guide
reserves $D$ (with $\partial$) for a scalar derivation of the surreal field,
and the collection's surcomplex-automorphisms report writes this same map
$S_a$ (\lbl{saut:eq:dilation}, there for every positive surreal $a$) while
using $D_\lambda$ and $D_\chi$ for character twists
$t^g\mapsto\chi(g)t^g$. Here $S_q$ always carries its subscript, and the
abbreviation $D$ is not used: every $D$, $D_2$, $D_q$, $D_r$, $D_s$ of the
source is $S$ with the same subscript ($S_2$ for the bare $D$). The only
$D$ left in the report is $D_\chi$ in Section~\ref{dsup:sec:position}, which
denotes that report's character twist. An unsubscripted $S$ is a well-ordered subset of $\Gamma$ (in the definition
of $K$ and in Lemma~\ref{dsup:lem:orbit}).
\item \emph{The support formula is $\Supp(m,x)$}, ``$m$ is a canonical
monomial in the support of $x$'' (equation~\eqref{dsup:eq:supportformula}).
The source wrote $S(m,x)$, which would now clash with the dilation. The
support \emph{set} is still $\supp(x)$.
\end{enumerate}

The other symbols are the source's. Their meanings here are fixed as
follows.
\begin{itemize}[leftmargin=*,itemsep=3pt]
\item $\sigma$ is always a coefficient-fixing monomial automorphism
$\sigma(t^g)=\chi(g)t^{\theta g}$ of the \emph{scalar} field
\eqref{dsup:eq:sigma}. It is not the argument dilation $f(z)\mapsto
f(\lambda z)$ of formal functions studied in the collection's
holonomic-rigidity report, and not the symbolic shift or the small-divisor
rate that the notation guide records under $\sigma$.
\item $\theta\in\Aut_{\mathrm{ord}}(\Gamma)$ is the surcomplex-automorphisms
report's $\tau$; $\chi$ and $\rho$ agree with its $\chi$ and $\rho$. In its
notation $\sigma=M_{\id,\theta,\chi}$ (\lbl{saut:thm:monomial}), and the maps
\eqref{dsup:eq:centralizerform} are its $M_{\rho,\theta}=M_{\rho,\theta,1}$.
\item $K=\kk((t^\Gamma))$ is a \emph{set-sized} full Hahn field over an
arbitrary coefficient field $\kk$ of characteristic zero; the guide's
$K_\Gamma=\C((t^\Gamma))$ is the case $\kk=\C$. (The surcomplex-automorphisms
report uses $K$ for the class $\No(i)$.) The class $\SC$ is the guide's
$\No[i]$. In Section~\ref{dsup:sec:surreal}, and in
Theorem~\ref{dsup:thm:surreal}(5) in particular, $\theta$ ranges over ordered
additive class automorphisms of $(\No,+,<)$: there the exponent group is a
proper class, departing from the guide's set-sized $\Gamma$ as the
surcomplex-automorphisms report also does.
\item $U_1=1+\mm$ is the group of principal units. It is not that report's
group $\mathcal U$ of $1$-automorphisms; the sets $\mathcal U_\pm(S)$ of
Lemma~\ref{dsup:lem:orbit} are orbit hulls of exponents. $M(m)$ is the
monomial formula~\eqref{dsup:eq:monomialformula}, not a map $M_{\rho,\tau}$.
$F$ is the fixed field of one $\sigma$ (Corollary~\ref{dsup:cor:fixed}) or of
a commuting family \eqref{dsup:eq:jointgroup}.
\item \emph{Resonance} means only the set
$R_\lambda=\{g:\theta g=g,\ \chi(g)=\lambda\}$ and the restriction
$P_\lambda$ to it \eqref{dsup:eq:resonance}. It is unrelated to the
multiplier resonances of the dynamics-and-normal-forms report.
\emph{Spectrum} is used once, in the algebraic sense stated after
Corollary~\ref{dsup:cor:fixed}. A \emph{centralizer} is always taken in a
group of field automorphisms, not among germs.
\item \emph{Strong} means preserving Hahn summation
(Definition in Section~\ref{dsup:sec:setup}), not convergence of finite
partial sums. \emph{Dilation} always means an exponent dilation of the
scalar field.
\end{itemize}

The tempting false readings, beside the true ones: $S_2$ is not squaring,
since $S_2(\sum a_gt^g)=\sum a_gt^{2g}$ while
$(\sum a_gt^g)^2=\sum a_ga_ht^{g+h}$, and the two agree exactly on the
canonical monomials (Lemma~\ref{dsup:lem:monomialrecognition}). $S_q$ is not
$f(z)\mapsto f(\lambda z)$, and it is not the surreal exponential, the
omega-map, a derivation or an exponential-preserving automorphism
(Section~\ref{dsup:sub:recovered}).

\subsection{Relation to other reports in the collection}\label{dsup:sec:position}

This subsection was written when the manuscript joined the collection; the
source compared itself only with the root README, the catalogue and the
surcomplex-automorphisms report's guide.

\paragraph{Surcomplex field automorphisms}
(\repo{surcomplex/surcomplex-field-automorphisms}). The dilation $S_q$ is
that report's $S_a$ at $a=q$ (\lbl{saut:eq:dilation}), and the maps
\eqref{dsup:eq:sigma} are its monomial automorphisms $M_{\rho,\tau,\chi}$
(\lbl{saut:thm:monomial}) with $\rho=\id$ and $\tau=\theta$. The source
credits that construction (Section~\ref{dsup:sec:intro}) and uses none of
that report's classification statements as premises. Four relations are
precise.
\begin{enumerate}[leftmargin=*,itemsep=3pt]
\item \emph{Consistency with the four-layer decomposition.}
\lbl{saut:thm:decomp} factors every valued automorphism $\alpha\in G_v$
(those with $\alpha(\OO)=\OO$) of $\No(i)$ uniquely as
$\alpha=u\,D_\chi M_{\rho,\tau}$. By Theorem~\ref{dsup:thm:surreal}(5), a
class field automorphism of $\No(i)$ commuting with some $S_q$, $q\ne1$, is
of the form $M_{\rho,\theta}$. So it lies in $G_v$, and in that
factorization $u=1$ and $\chi=1$. Theorem~\ref{dsup:thm:centralizer} is the
set-sized analogue, and the computation after
Corollary~\ref{dsup:cor:nohidden} is the same statement for the layer
$D_\chi$ alone ($D_\chi$ commutes with $S_2$ only if $\chi=1$). Both
uniqueness statements describe the same maps; there is no conflict.
\item \emph{One subgroup beyond $G_v$.} That report states, as non-claims,
that its four-layer decomposition covers $G_v$ only, not all field
automorphisms, and that no complete classification of the automorphisms of
$\No(i)$ is claimed (items~1 and~10 of \lbl{saut:sec:nonclaims}). The
centralizer here is computed among \emph{all} field automorphisms, with no
assumption that they preserve the valuation ring or Hahn sums; being valued
and strong is a conclusion. Thus one subgroup of the unrestricted group,
the centralizer of $S_q$, is shown to lie in $G_v$ and is identified there.
That is the whole extension. It does not classify $\Aut(\No(i))$, it says
nothing about automorphisms that commute with no $S_q$, and both non-claims
of that report otherwise stand.
\item \emph{Complement to non-definability.}
\lbl{saut:prop:nondefinability} shows that over every set of parameters
there are pure-field automorphisms moving $\No$ and moving $\OO$, so neither
is definable in the pure field with parameters.
Theorem~\ref{dsup:thm:definability} and
Theorem~\ref{dsup:thm:surreal}(4) show that after adding the single
automorphism $S_q$ the coefficient field, the monomials, every coefficient
and the valuation ring become parameter-free definable. The expansion is a
different structure, so the two results are complementary. On $\No(i)$ the
real axis $\No$ is still not claimed definable, and the centralizer
contains every automorphism of $\C$ (Section~\ref{dsup:sub:recovered}).
\item \emph{Its questions.} None of the four questions of
\lbl{saut:sec:questions} is answered here: the image of exponential
automorphisms on the value group, omega-map compatibility, effective
descriptions inside the leading-term kernel, and real-form and continuity
classifications. The source says that it does not settle the
exponential-automorphism image problems discussed in \cite{KKS,RepoAut}
(Section~\ref{dsup:sub:recovered}). The only centralizer computed in that
report is that of conjugation (\lbl{saut:thm:axis}), a different group; and
that report never asserts that $S_r$ and $S_s$ are conjugate, so
Theorem~\ref{dsup:thm:nonconjugate} contradicts nothing there.
\end{enumerate}

\paragraph{Exponential automorphism rigidity}
(\repo{surreal/exponential-automorphism-rigidity}). On $\No$, $S_q$ is that
report's canonical lift $\widehat T_q$ (\lbl{prop:hahn-lift}) of the
value-group dilation $T_q(\gamma)=q\gamma$, and for rational $q\ne1$ that
report proves that $T_q$ has no exponential lift (\lbl{thm:dilation}). With
$\kk=\R$, Theorem~\ref{dsup:thm:surreal}(5) says that the class field
automorphisms of $\No$ commuting with $S_q$ are exactly the canonical lifts
$\widehat T$ of the ordered additive automorphisms $T$ of $(\No,+,<)$.
Theorems~\ref{dsup:thm:rationaldef} and~\ref{dsup:thm:nonconjugate} are new
statements about the same maps and contradict nothing there. Nothing here
decides which of these lifts, or which other automorphisms, commute with
the exponential; that report's Question~13.1 (\lbl{q:image}) is untouched.

\paragraph{Omnific-preserving automorphisms}
(\repo{surreal/omnific-preserving-automorphisms}). \emph{Related (batch 34).}
Part~X of that report, from a manuscript that does not cite this one,
re-derives the difference core of this report with trivial character. Its
weighted Green operator (\lbl{opa:us:thm:green}) is the partial inverse of
Theorem~\ref{dsup:thm:resolvent} with $\chi=1$, for the monomial automorphism
of any ordered additive automorphism of a set or class exponent group; the
dictionary is \lbl{opa:us:rem:dsup}. Its proof works directly on the class
$\No$ instead of localizing to a set-sized subgroup. For the automorphisms of
its faithful actions of left-orderable groups on $\No$ and $\SC$, its exact
multiplicative obstruction (\lbl{opa:us:thm:multiplicative}) and its relative
multiplicative retraction (\lbl{opa:us:thm:retraction}) specialize
Theorem~\ref{dsup:thm:multsingle}; the retraction adds an explicit image of
the inner difference operator. Its variable-coefficient criterion
(\lbl{opa:us:thm:firstorder}) solves $\sigma x-ax=b$ for nonconstant $a$,
beyond the constant-scalar resolvent here. Its central cocycle normal forms
(\lbl{opa:us:thm:cocycles}) give the degree-one cohomology for every group
with a central nonidentity element, where
Theorem~\ref{dsup:thm:dilationcohom} computes every degree for $\Z^d$ acting
by rational dilations. No statement of this report changes.

\paragraph{Holonomic rigidity}
(\repo{surcomplex/holonomic-rigidity-for-entire-hahn-functions}). The
dilations in that report's batch-22 material are argument dilations
$f(z)\mapsto f(\lambda z)$ of formal functions, with $\lambda$ in the field.
The dilations here act on the exponents of scalars. The two share a word,
not a theorem.

\paragraph{Birthday cutoffs and hereditary sets}
(\repo{foundations-and-computation/birthday-cutoffs-and-hereditary-sets}),
placed after the pin. There the ordered surreal field expanded by the
birthday function interprets full second-order arithmetic
(\lbl{hset:prop:arithmetic}), its complete theory is undecidable
(\lbl{hset:cor:undecidable}), and a definable relation has the independence
property (\lbl{hset:prop:ip}). Section~\ref{dsup:sec:logic} obtains the same
kind of consequence, and in addition $\mathrm{TP}_2$ and the strict order
property, for a different expansion (a Hahn field with one exponent
dilation) by a different mechanism (Boolean coefficient coding, after
Camacho \cite{Camacho}). Neither result implies the other.

\paragraph{Unrelated uses of shared words.} The resonances, homological
operator and centralizers of \repo{surcomplex/dynamics-and-normal-forms}
concern local dynamics of germs, and the Koszul complexes of
\repo{surcomplex/analytic-geometry} and \repo{surcomplex/finite-deformations}
are complexes of polynomial ideals. No statement there is used or extended.
No \texttt{dsup:} label has a Lean mapping in \repo{docs/FORMALIZATION.md}.

\section{Conventions, supports, and monomial automorphisms}\label{dsup:sec:setup}

Throughout the set-sized theory, $\Gamma$ is an ordered abelian group,
$\kk$ is a field of characteristic zero, and
\[
 K=\kk((t^\Gamma))
 =\left\{\sum_{g\in S}a_gt^g:S\subseteq\Gamma
       \text{ is well ordered},\ a_g\in\kk\right\}
\]
is the \emph{full} Hahn field. Zero coefficients can be omitted. We write
$x_g$ for the coefficient at $g$, $\supp(x)=\{g:x_g\ne0\}$, and
\[
 v(x)=\min\supp(x)\quad(x\ne0),\qquad v(0)=\infty.
\]
The coefficient field is embedded by $a\mapsto at^0$. The valuation ring,
maximal ideal, and principal-unit group are
\[
 \OO=\{x:v(x)\ge0\},\qquad
 \mm=\{x:v(x)>0\},\qquad U_1=1+\mm.
\]
The notation $[x]_0=x_0$ always means the coefficient of $t^0$, including
for $x\notin\OO$. It is not a ring homomorphism on all of $K$.

\begin{definition}[Hahn summability]
A set-indexed family $(x_i)_{i\in I}$ is \emph{summable} if
$\bigcup_i\supp(x_i)$ is well ordered and, for each $g$, only finitely many
$i$ have $(x_i)_g\ne0$. Its sum has coefficient
$\sum_i(x_i)_g$. A $\kk$-linear map is \emph{strong} if it takes summable
families to summable families and commutes with their sums.
\end{definition}

Restriction to any fixed subset of $\Gamma$ is a strong linear map, since
it only deletes support points. Strong summation permits infinite families
whose finite subsums do not converge in the valuation topology. In
particular, no argument below replaces a support proof by the assertion
that an orbit ``tends to infinity.''

\begin{definition}[Coefficient-fixing monomial automorphism]
Let $\theta\in\Aut_{\mathrm{ord}}(\Gamma)$ and let
$\chi:\Gamma\to\kk^\times$ be a group homomorphism. Put
\begin{equation}\label{dsup:eq:sigma}
 \sigma\left(\sum_g a_gt^g\right)
   =\sum_g a_g\chi(g)t^{\theta g}.
\end{equation}
\end{definition}

\begin{proposition}\label{dsup:prop:monomial}
The map \eqref{dsup:eq:sigma} is a strong field automorphism fixing $\kk$.
Its valuation action is $v(\sigma x)=\theta(vx)$ for $x\ne0$. Two such
maps $\sigma_i,\sigma_j$ commute if and only if
\begin{equation}\label{dsup:eq:commute-data}
 \theta_i\theta_j=\theta_j\theta_i,\qquad
 \chi_i(\theta_jg)\chi_j(g)=\chi_j(\theta_i g)\chi_i(g)
 \quad(g\in\Gamma).
\end{equation}
\end{proposition}
\begin{proof}
An increasing bijection preserves well-ordered supports and point-finiteness.
Additivity and strength follow coefficientwise. On monomials,
$\chi(g+h)=\chi(g)\chi(h)$ and $\theta(g+h)=\theta g+\theta h$ give
multiplicativity. Hahn convolution is legitimate because its coefficient
sums are finite; thus the monomial calculation extends to arbitrary
series. The inverse sends $t^g$ to
$\chi(\theta^{-1}g)^{-1}t^{\theta^{-1}g}$. Leading terms give the valuation
formula. Computing both compositions on $t^g$ proves
\eqref{dsup:eq:commute-data}, and strength extends the result to all series.
\end{proof}

For $q\in\Q_{>0}$ on a divisible $\Gamma$, write $S_q$ for the case
$\theta g=qg$ and $\chi=1$. For $S_2$ alone, $2$-divisibility is sufficient.
No ordinary absolute value on $\kk$ is used in the Hahn theory.

\section{The orbit-support lemma}\label{dsup:sec:orbits}

The following elementary support fact is responsible for the unrestricted
rank of the results. The forward direction of summation is chosen
separately on the two movement chambers of $\theta$.

\begin{lemma}[Increasing orbit hull]\label{dsup:lem:orbit}
Let $\theta$ be an increasing automorphism of an ordered abelian group.
If $S\subseteq\{g:\theta g>g\}$ is well ordered, then
\[
 \mathcal U_+(S)=\{\theta^n s:s\in S,\ n\in\N\}
\]
is well ordered, and every exponent has only finitely many representations
$\theta^n s$ with $(n,s)\in\N\times S$.
If $S\subseteq\{g:\theta g<g\}$ is well ordered, the same assertions hold
for
$\mathcal U_-(S)=\{\theta^{-n}s:s\in S,\ n\ge1\}$.
\end{lemma}
\begin{proof}
Every sequence in a well-ordered set has an infinite nondecreasing
subsequence. One proof uses tail minima: if a tail minimum occurs
infinitely often, take those occurrences; otherwise pass beyond its last
occurrence and continue with the larger minimum of the remaining tail.
If this process does not stop, the chosen values strictly increase.

Suppose a strictly decreasing sequence in $\mathcal U_+(S)$ existed.
Choose representations $\theta^{n_j}s_j$. By two successive subsequence
selections, both $n_j$ and $s_j$ may be taken nondecreasing. For $j<k$,
\[
 \theta^{n_j}s_j\le\theta^{n_j}s_k\le\theta^{n_k}s_k,
\]
because the forward orbit of $s_k$ increases. This contradicts strict
decrease. In ZFC, a linear order with no infinite strictly decreasing
sequence is well ordered: a nonempty subset without a least element
would recursively supply such a sequence.

If one exponent $g$ had infinitely many representations, their $n$'s
would be unbounded, since each $n$ determines at most one $s$. Passing to
strictly increasing $n_j$ gives
$s_j=\theta^{-n_j}g$. The point $g$ is in the increasing chamber, which
is $\theta$-invariant, so these $s_j$ strictly decrease. That contradicts
the well-ordering of $S$.

For the second assertion apply the first to $\theta^{-1}$ on the decreasing
chamber and omit the $n=0$ terms. Every subset of a well-ordered set is
well ordered.
\end{proof}

\begin{corollary}[Labeled version]\label{dsup:cor:labeled}
Suppose $(x_i)_{i\in I}$ is a summable family whose supports are all in the
increasing chamber. Then $(\sigma^n x_i)_{(n,i)\in\N\times I}$ is summable.
The corresponding inverse-iterate assertion holds on the decreasing
chamber.
\end{corollary}
\begin{proof}
Apply Lemma~\ref{dsup:lem:orbit} to the union of the original supports. For a
given output exponent there are finitely many pairs $(n,s)$; at each $s$
there are finitely many original labels $i$. Nonzero character factors
and nonzero scalar multipliers do not affect this support argument.
\end{proof}

\begin{remark}
The lemma is not an assertion of cofinal growth. For example, in
$\Q^2$ with lexicographic order, the shear
$\theta(a,b)=(a,b+a)$ takes $(1,0)$ through $(1,n)$, a sequence bounded
above by $(2,0)$. Its orbit is summable in the Hahn sense, but finite
subsums need not converge in the valuation topology.
\end{remark}

\section{An exact resolvent and its resonance obstruction}\label{dsup:sec:resolvent}

Fix $\sigma$ as in \eqref{dsup:eq:sigma} and $\lambda\in\kk^\times$. Define
\[
 A_+=\{g:\theta g>g\},\quad A_-=\{g:\theta g<g\},\quad
 A_=\Fix(\theta).
\]
Write $f=f_++f_-+f_=$ by restricting its support to these three sets. The
\emph{resonance set} and \emph{resonance projection} are
\begin{equation}\label{dsup:eq:resonance}
 R_\lambda=\{g\in\Gamma:\theta g=g,\ \chi(g)=\lambda\},
 \qquad P_\lambda f=\sum_{g\in R_\lambda}f_gt^g.
\end{equation}
Define
\begin{align}
 Q_\lambda f={}&-\sum_{n\ge0}\lambda^{-n-1}\sigma^n(f_+)
        +\sum_{n\ge1}\lambda^{n-1}\sigma^{-n}(f_-)
 \notag\\[-2pt]
 &+\sum_{\substack{g\in A_=\\\chi(g)\ne\lambda}}
           \frac{f_g}{\chi(g)-\lambda}t^g .\label{dsup:eq:Q}
\end{align}
All sums here are Hahn sums.

\begin{theorem}[Strong partial inverse]\label{dsup:thm:resolvent}
The map $Q_\lambda:K\to K$ is well defined, strong, and $\kk$-linear.
It satisfies
\begin{align}
 (\sigma-\lambda)Q_\lambda
   &=Q_\lambda(\sigma-\lambda)=I-P_\lambda,\label{dsup:eq:resolvent-identities}\\
 P_\lambda Q_\lambda&=Q_\lambda P_\lambda=0,\label{dsup:eq:normalization}\\
 v(Q_\lambda f)&\ge v(f).\label{dsup:eq:valuation-Q}
\end{align}
Consequently,
\begin{equation}\label{dsup:eq:kerim}
 \ker(\sigma-\lambda)=P_\lambda K,
 \qquad \im(\sigma-\lambda)=\ker P_\lambda.
\end{equation}
The equation $(\sigma-\lambda)y=f$ is solvable exactly when
$P_\lambda f=0$. Its unique solution with $P_\lambda y=0$ is $Q_\lambda f$.
\end{theorem}
\begin{proof}
Lemma~\ref{dsup:lem:orbit} makes each of the first two families in \eqref{dsup:eq:Q}
summable. The third is a restriction and coefficient rescaling of $f$.
The finite union of the three supports is well ordered and the three
families remain point-finite, so their sum exists. The labeled version
proves strength and justifies interchanging any summable input family
with the displayed sums.

On the increasing chamber, apply $\sigma-\lambda$ to the first sum:
\[
 -\sum_{n\ge0}\lambda^{-n-1}\sigma^{n+1}f_+
 +\sum_{n\ge0}\lambda^{-n}\sigma^n f_+=f_+.
\]
The cancellation is coefficientwise between jointly summable families.
There is no limiting remainder. On the decreasing chamber the calculation is
\[
 \sum_{n\ge1}\lambda^{n-1}\sigma^{-(n-1)}f_-
 -\sum_{n\ge1}\lambda^n\sigma^{-n}f_-=f_-.
\]
On a fixed exponent, $\sigma-\lambda$ multiplies its coefficient by
$\chi(g)-\lambda$. This proves the left-hand identity in
\eqref{dsup:eq:resolvent-identities}. The chamber restrictions commute with
$\sigma$, as do the fixed-exponent division operation and the iterate
sums. Thus $Q_\lambda$ commutes with $\sigma$, proving the right-hand
identity. No term of \eqref{dsup:eq:Q} lies in $R_\lambda$, which gives
\eqref{dsup:eq:normalization}.

Every exponent appearing in the first sum is at least its originating
exponent, because it is a forward iterate in the increasing chamber.
The same is true of inverse iterates in the decreasing chamber. Fixed
exponents are unchanged. This proves \eqref{dsup:eq:valuation-Q}, including the
possibility of cancellation.

For completeness, the kernel also has a direct support proof. If
$\sigma y=\lambda y$, then $\theta\supp(y)=\supp(y)$. A nonfixed exponent
would therefore force its whole bi-infinite orbit into the support. That
orbit contains an infinite strictly decreasing sequence. Hence the support
is fixed pointwise by $\theta$, and its coefficients require
$\chi(g)=\lambda$. This gives the kernel in \eqref{dsup:eq:kerim}.
The projection $P_\lambda$ annihilates the image of $\sigma-\lambda$;
conversely \eqref{dsup:eq:resolvent-identities} produces every element of its
kernel. Finally, two normalized solutions differ by a kernel element with
zero resonance projection and therefore coincide.
\end{proof}

\begin{corollary}[Fixed fields and eigenspaces]\label{dsup:cor:fixed}
Let
\[
 \Gamma_\sigma=\{g:\theta g=g,\ \chi(g)=1\},\qquad
 F=\kk((t^{\Gamma_\sigma})).
\]
Then $K^\sigma=F$. If $R_\lambda$ is nonempty and $g_\lambda\in R_\lambda$,
then
\[
 R_\lambda=g_\lambda+\Gamma_\sigma,\qquad
 \ker(\sigma-\lambda)=t^{g_\lambda}F.
\]
For every $n\ge1$,
$\ker(\sigma-\lambda)^n=\ker(\sigma-\lambda)$.
\end{corollary}
\begin{proof}
The set $\Gamma_\sigma$ is a subgroup, so its full Hahn series form a
subfield. The fixed-field statement is Theorem~\ref{dsup:thm:resolvent} with
$\lambda=1$. Subtracting $g_\lambda$ from another resonant exponent proves
the coset formula, and translation of supports gives the eigenspace.
The decomposition $K=P_\lambda K\oplus\ker P_\lambda$ is invariant under
$\sigma-\lambda$, which is zero on its first summand and invertible on
its second. Powers have the same kernel.
\end{proof}

This is a statement about kernels and invertibility, not a claim that
$K$ is a direct sum of its eigenspaces. In particular, arbitrary series
can have support on moving orbits. In the purely algebraic sense in which
the spectrum of a linear map consists of scalars where it fails to be
bijective, the spectrum of $\sigma$ is
$\chi(\Fix\theta)$. Zero is not in this spectrum.

\begin{corollary}[Split polynomial difference operators]\label{dsup:cor:poly}
Suppose $p\in\kk[X]\setminus\{0\}$ splits over $\kk$. Put
\[
 R_p=\{g:\theta g=g,\ p(\chi(g))=0\},\qquad P_pf=\restr f{R_p}.
\]
Then $\ker p(\sigma)=P_pK$ and $\im p(\sigma)=\ker P_p$.
Repeated roots of $p$ do not enlarge the kernel. On $\ker P_p$, the inverse
is strong.
\end{corollary}
\begin{proof}
The factor $X$ is harmless since $\sigma$ is a strong automorphism.
On the moving chambers every nonzero scalar factor $\sigma-\lambda$ has
the inverse supplied by \eqref{dsup:eq:Q}; these inverses commute because they
are inverses of commuting bijections on that invariant subspace.
On the fixed chamber $p(\sigma)$ multiplies the coefficient at $g$ by
$p(\chi(g))$. Invert this scalar where nonzero and set the inverse to zero
on $R_p$. The factors and all restrictions preserve summable families.
\end{proof}

\begin{remark}[Not a residue homomorphism]
For $\lambda=1$, $P_1$ is $F$-linear: multiplication by a fixed monomial
translates supports by an element of $\Gamma_\sigma$, and does not change
the resonance test; general $F$-linearity follows by Hahn convolution.
Nevertheless, $P_1$ need not be multiplicative on $K$. For doubling and
$g\ne0$, both $P_1(t^g)$ and $P_1(t^{-g})$ vanish, whereas
$P_1(t^gt^{-g})=1$.
\end{remark}

\section{Commuting equations and an explicit Koszul contraction}\label{dsup:sec:koszul}

Let $\sigma_1,\ldots,\sigma_d$ be a finite commuting family of
coefficient-fixing monomial automorphisms. Write
\[
 \Delta_j=\sigma_j-I,\quad P_j=P_{j,1},\quad Q_j=Q_{j,1},\quad
 E_0=I,\quad E_j=P_1\cdots P_j,\quad P=E_d.
\]
The subscript $(j,1)$ refers to the construction in
Theorem~\ref{dsup:thm:resolvent} for $\sigma_j$ and $\lambda=1$. Let
\begin{equation}\label{dsup:eq:jointgroup}
 \Gamma_0=\{g:\theta_jg=g,\ \chi_j(g)=1\text{ for every }j\},
 \qquad F=\kk((t^{\Gamma_0})).
\end{equation}
Thus $P$ restricts a series to its joint-resonant exponents.

\begin{lemma}[Commuting normalized inverses]\label{dsup:lem:operatorscommute}
The maps $P_i,Q_i,\sigma_i$ commute with $P_j,Q_j,\sigma_j$ whenever the
two displayed maps have different indices. In particular $P$ commutes
with every operator, and $F=K^{\sigma_1,\ldots,\sigma_d}$.
\end{lemma}
\begin{proof}
If $g$ is fixed by $\theta_i$, commutativity of the exponent maps makes
$\theta_jg$ fixed by $\theta_i$. The character identity
\eqref{dsup:eq:commute-data} then gives
$\chi_i(\theta_jg)=\chi_i(g)$. Thus the resonance subset for $i$ is
invariant under $\theta_j$ and its inverse. Restriction to this subset
commutes with $\sigma_j$. The projections also commute with one another,
being support restrictions.

Here is a useful uniqueness argument for the inverse. If a linear map
$A$ commutes with $\Delta_i$ and $P_i$, then $AQ_i f$ and $Q_iAf$ both
have zero $P_i$-projection and both solve
\[
 \Delta_i z=(I-P_i)Af.
\]
The normalized solution is unique, so $AQ_i=Q_iA$. Apply this first to
$A=\sigma_j$ and $A=P_j$. It follows that $Q_i$ commutes with $\Delta_j$
and $P_j$; applying the same argument for $Q_j$ then gives $Q_iQ_j=Q_jQ_i$.
The common fixed field is obtained by intersecting the fixed-support
conditions of Corollary~\ref{dsup:cor:fixed}.
\end{proof}

Let $e_1,\ldots,e_d$ be the standard basis of $\kk^d$, and set
\[
 C^r=K\otimes_\kk\bigwedge^r\kk^d,
 \qquad
 \dd=\sum_{j=1}^d\Delta_j\,\eps_j,
 \qquad \eps_j(\eta)=e_j\wedge\eta.
\]
Here a cochain is only a finite tuple of Hahn series. The usual contraction
$\iota_j$ on the exterior algebra is defined by
\[
 \iota_j(e_{i_1}\wedge\cdots\wedge e_{i_r})
 =\begin{cases}
 (-1)^{s-1}e_{i_1}\wedge\cdots\widehat{e_{i_s}}\cdots\wedge e_{i_r},
        &i_s=j,\\
 0,&j\notin\{i_1,\ldots,i_r\}.
 \end{cases}
\]
It satisfies
$\eps_i\iota_j+\iota_j\eps_i=\delta_{ij}I$.
Extend $P$ to each cochain degree by applying it to every coefficient;
call this extension $\Pi$.

\begin{theorem}[Explicit deformation retraction]\label{dsup:thm:koszul}
The operator
\begin{equation}\label{dsup:eq:homotopy}
 \hh=\sum_{j=1}^d Q_jE_{j-1}\iota_j:C^r\longrightarrow C^{r-1}
\end{equation}
is strong on coefficient tuples and satisfies
\begin{equation}\label{dsup:eq:homotopyidentity}
 \dd\hh+\hh\dd=I-\Pi,\qquad
 \hh^2=0,\qquad \hh\Pi=\Pi\hh=0,\qquad
 \dd\Pi=\Pi\dd=0.
\end{equation}
Consequently
\begin{equation}\label{dsup:eq:cohomology}
 H^r(\Z^d,K_{\mathrm{add}})
 \simeq F\otimes_\kk\bigwedge^r\kk^d
 \simeq F^{\binom dr}
 \quad(0\le r\le d),
\end{equation}
and the cohomology vanishes for $r>d$. The isomorphism sends a cocycle to
its coefficientwise joint-resonant part.
\end{theorem}
\begin{proof}
All coefficient operators in \eqref{dsup:eq:homotopy} commute by
Lemma~\ref{dsup:lem:operatorscommute}. Therefore the exterior-algebra identities
reduce the anticommutator to its diagonal terms:
\begin{align*}
 \dd\hh+\hh\dd
 &=\sum_{j=1}^d\Delta_jQ_jE_{j-1}\\
 &=\sum_{j=1}^d(I-P_j)E_{j-1}
   =\sum_{j=1}^d(E_{j-1}-E_j)=I-P.
\end{align*}
The contractions anticommute and square to zero, so $\hh^2=0$. Each term
$Q_jE_{j-1}$ annihilates $P$ because $Q_jP_j=0$; commutativity gives both
$\hh\Pi=0$ and $\Pi\hh=0$. Finally $\Delta_jP=P\Delta_j=0$ for every $j$,
which proves the remaining identities. Strength follows from strength of
the finitely many factors and sums.

The target subcomplex $F\otimes\bigwedge^\bullet\kk^d$ has zero
differential. If $\dd\alpha=0$, then
$\alpha=\dd\hh\alpha+\Pi\alpha$, and a coboundary has zero projection.
Thus the inclusion of that zero-differential complex induces exactly
\eqref{dsup:eq:cohomology}.

To identify this with group cohomology, regard $K_{\mathrm{add}}$ as a
module over $\Z[z_1^{\pm1},\ldots,z_d^{\pm1}]$, with $z_j$ acting by
$\sigma_j$. The tensor product of the two-term free resolutions
\[
 0\longrightarrow\Z[z_j^{\pm1}]
   \xrightarrow{\ z_j-1\ }\Z[z_j^{\pm1}]
   \longrightarrow\Z\longrightarrow0
\]
is a free resolution of the trivial module for $\Z^d$. Exactness follows
successively because, after imposing earlier relations, the next
$z_j-1$ remains a non-zero-divisor in a Laurent polynomial ring.
Applying $\operatorname{Hom}$ into $K_{\mathrm{add}}$ produces the complex
above, up to the standard consistent choice of exterior signs. It has
length $d$, giving the assertion in larger degrees.
\end{proof}

\begin{corollary}[Simultaneous additive equations]\label{dsup:cor:simultaneous}
Let $f_1,\ldots,f_d\in K$. The system
\[
 \Delta_jy=f_j\qquad(1\le j\le d)
\]
has a solution if and only if
\begin{equation}\label{dsup:eq:additivecriteria}
 \Delta_i f_j=\Delta_jf_i\quad\text{for all }i,j,
 \qquad Pf_j=0\quad\text{for all }j.
\end{equation}
When these conditions hold, its unique solution with $Py=0$ is
\begin{equation}\label{dsup:eq:simultaneoussolution}
 y=\sum_{j=1}^d Q_jE_{j-1}f_j.
\end{equation}
All other solutions are obtained by adding an arbitrary element of $F$.
\end{corollary}
\begin{proof}
The first condition says that $\alpha=\sum_j f_je_j$ is a cocycle. The
second says $\Pi\alpha=0$. Equation~\eqref{dsup:eq:homotopyidentity} gives
$\dd(\hh\alpha)=\alpha$, and contraction of the basis vectors gives
\eqref{dsup:eq:simultaneoussolution}. Necessity follows from
$\dd^2=0$ and $\Pi\dd=0$. The normalization follows from $P\hh=0$, and
uniqueness from the common fixed-field calculation.
\end{proof}

\begin{corollary}[Definable invariant projection]\label{dsup:cor:defproj}
In the language of fields with $\sigma_1,\ldots,\sigma_d$, the graph of
$P:K\to F$ is parameter-free definable by
\begin{equation}\label{dsup:eq:defproj}
 P(x)=c\quad\Longleftrightarrow\quad
 \bigwedge_{j=1}^d\sigma_jc=c
 \ \land\ 
 \exists y_1\cdots y_d\left(x-c=\sum_{j=1}^d\Delta_jy_j\right).
\end{equation}
\end{corollary}
\begin{proof}
The operator identity
$I-P=\sum_j\Delta_jQ_jE_{j-1}$ shows that the right-hand side holds for
$c=Px$. Conversely $P$ kills every $\Delta_jy_j$ and fixes $c\in F$, so
applying $P$ to the displayed equality gives $Px=c$.
\end{proof}

\begin{remark}[A cohomology calculation, not a false Leibniz rule]
The operator $\Delta_j$ satisfies a twisted product rule, not the ordinary
Leibniz rule. We use the Koszul complex as a complex of additive groups.
No claim is made that its naive coefficientwise wedge multiplication
makes $\dd$ an ordinary derivation. This avoids confusing its additive
calculation with a differential graded algebra structure.
\end{remark}

\section{Principal units and multiplicative difference equations}\label{dsup:sec:multiplicative}

\subsection{The formal logarithm is the appropriate logarithm}

For $h\in\mm$, the formal series
\begin{equation}\label{dsup:eq:explog}
 \exp(h)=\sum_{n\ge0}\frac{h^n}{n!},\qquad
 \log(1+h)=\sum_{n\ge1}\frac{(-1)^{n+1}}n h^n
\end{equation}
are strongly summable. The underlying support fact is the Neumann lemma:
the additive monoid generated by a well-ordered subset of $\Gamma_{>0}$
is well ordered, with only finitely many contributions to a fixed
exponent across all finite word lengths. This standard Hahn-series
foundation, also part of the repository's documented summation layer
\cite{Repo}, is used here in precisely its formal-support sense.

The formal power-series identities for exponential and logarithm can
therefore be evaluated in $\mm$. Joint summability justifies substitution
and regrouping. They give mutually inverse group isomorphisms
\[
 (\mm,+)\ \underset{\log}{\stackrel{\exp}{\rightleftarrows}}\ (U_1,\cdot).
\]
A strong coefficient-fixing automorphism commutes with both maps, since
it preserves the rational coefficients and every summable family in
\eqref{dsup:eq:explog}. This is an infinitesimal Hahn exponential. It is not a
claim that a monomial automorphism commutes with the global surreal
exponential.

\begin{lemma}[Multiplicative coordinates]\label{dsup:lem:multcoords}
Every $x\in K^\times$ has a unique expression
\begin{equation}\label{dsup:eq:multcoords}
 x=c\,t^g\exp(h),\qquad
 g=v(x),\quad c=\lc(x)\in\kk^\times,\quad h\in\mm.
\end{equation}
Under $\sigma_j$, these coordinates transform as
\begin{equation}\label{dsup:eq:multaction}
 (g,c,h)\longmapsto(\theta_jg,\ c\chi_j(g),\ \sigma_jh).
\end{equation}
\end{lemma}
\begin{proof}
Dividing $x$ by its leading monomial $ct^g$ leaves a principal unit; apply
its unique logarithm. Multiplication adds the exponent and logarithm
coordinates and multiplies the coefficient coordinate. Applying
\eqref{dsup:eq:sigma} and the naturality of \eqref{dsup:eq:explog} gives
\eqref{dsup:eq:multaction}.
\end{proof}

Put $\mm_F=F\cap\mm$. The projections $P_j$ preserve $\mm$, and
\eqref{dsup:eq:valuation-Q} shows that $Q_j$ also preserves it. Thus the whole
Koszul contraction restricts to infinitesimals.

\begin{corollary}[Principal-unit cohomology]\label{dsup:cor:unitcohom}
For the commuting family of Section~\ref{dsup:sec:koszul},
\[
 H^r(\Z^d,\mm)\simeq\mm_F^{\binom dr},\qquad
 H^r(\Z^d,U_1)\simeq\mm_F^{\binom dr}
 \quad(0\le r\le d),
\]
where the second isomorphism uses $\log$. Both vanish above degree $d$.
\end{corollary}

\subsection{One multiplicative equation}

\begin{theorem}[Exact multiplicative obstruction]\label{dsup:thm:multsingle}
For one automorphism $\sigma$, write
$f=c\,t^\gamma\exp(h)$ in multiplicative coordinates and let $P=P_1$,
$Q=Q_1$. The equation
\begin{equation}\label{dsup:eq:multsingle}
 \frac{\sigma y}{y}=f,\qquad y\in K^\times,
\end{equation}
is solvable if and only if there exists $g\in\Gamma$ with
\begin{equation}\label{dsup:eq:multsinglecriteria}
 (\theta-I)g=\gamma,\qquad \chi(g)=c,
 \qquad Ph=0.
\end{equation}
For each such $g$, $y=t^g\exp(Qh)$ is a solution. The full solution set,
when nonempty, is a torsor under $(K^\sigma)^\times$.
\end{theorem}
\begin{proof}
Write $y=d\,t^g\exp(u)$. Formula~\eqref{dsup:eq:multaction} gives
\begin{equation}\label{dsup:eq:multquotient}
 \frac{\sigma y}{y}
 =\chi(g)t^{(\theta-I)g}\exp((\sigma-I)u).
\end{equation}
Uniqueness of multiplicative coordinates proves necessity of the first
two conditions and reduces the third to $\Delta u=h$. Theorem
\ref{dsup:thm:resolvent} solves that equation exactly when $Ph=0$, with
$u=Qh\in\mm$. The ratio of any two solutions is fixed by $\sigma$, and
multiplication by a fixed nonzero element preserves a solution.
\end{proof}

For a quotient description, use componentwise addition/multiplication on
$\Gamma\times\kk^\times$. The map
\[
 g\longmapsto((\theta-I)g,\chi(g))
\]
is a group homomorphism. The preceding theorem yields
\begin{equation}\label{dsup:eq:multquotientgroup}
 K^\times/\{\sigma y/y:y\in K^\times\}
 \simeq
 \frac{\Gamma\times\kk^\times}
      {\{((\theta-I)g,\chi(g)):g\in\Gamma\}}
 \ \times\ \mm_{K^\sigma}.
\end{equation}
The isomorphism sends the class of $c t^\gamma\exp(h)$ to the class of
$(\gamma,c)$ and to $Ph$. Its kernel is exactly the displayed subgroup.
In general the first quotient cannot be replaced by independent
``valuation'' and ``leading coefficient'' quotients: the same $g$ must
solve both equations in \eqref{dsup:eq:multsinglecriteria}.

\subsection{A commuting multiplicative system}

A tuple $(f_1,\ldots,f_d)$ in $(K^\times)^d$ is multiplicatively
compatible if
\begin{equation}\label{dsup:eq:multcompatible}
 \frac{\sigma_i(f_j)}{f_j}=\frac{\sigma_j(f_i)}{f_i}
 \quad\text{for all }i,j.
\end{equation}
This is precisely the consistency condition forced by
$\sigma_jy/y=f_j$.

\begin{theorem}[Simultaneous multiplicative equations]\label{dsup:thm:multsystem}
Write $f_j=c_jt^{\gamma_j}\exp(h_j)$. Suppose
\eqref{dsup:eq:multcompatible} holds. Then there exists $y\in K^\times$ with
$\sigma_jy/y=f_j$ for every $j$ if and only if
\begin{equation}\label{dsup:eq:multsystemcriteria}
 \begin{gathered}
 \text{there is one }g\in\Gamma\text{ such that }
    (\theta_j-I)g=\gamma_j,\quad \chi_j(g)=c_j\quad(1\le j\le d),\\
 Ph_j=0\quad(1\le j\le d).
 \end{gathered}
\end{equation}
When these conditions hold, a solution is
\begin{equation}\label{dsup:eq:multsystemsolution}
 y=t^g\exp\left(\sum_{j=1}^d Q_jE_{j-1}h_j\right),
\end{equation}
and all solutions differ by an element of $F^\times$.
\end{theorem}
\begin{proof}
Applying uniqueness of multiplicative coordinates to
\eqref{dsup:eq:multcompatible} gives exactly
\begin{align*}
 (\theta_i-I)\gamma_j&=(\theta_j-I)\gamma_i,\\
 \chi_i(\gamma_j)&=\chi_j(\gamma_i),\\
 \Delta_i h_j&=\Delta_j h_i.
\end{align*}
The constants $c_j$ cancel from compatibility, but not from the search
for a common solution. Formula~\eqref{dsup:eq:multquotient}, for all $j$, shows
that the exponent and coefficient coordinates require a single $g$ as
in \eqref{dsup:eq:multsystemcriteria}. The remaining additive system in
$\mm$ is solved by Corollary~\ref{dsup:cor:simultaneous} and its infinitesimal
restriction. This gives \eqref{dsup:eq:multsystemsolution}. Ratios of solutions
are fixed by every $\sigma_j$.
\end{proof}

For clarity about the residual cohomology, let
$M=\Gamma\times\kk^\times$ with action
$(g,c)\mapsto(\theta_jg,\chi_j(g)c)$. Lemma~\ref{dsup:lem:multcoords} is an
equivariant isomorphism of abelian groups $K^\times\simeq M\times\mm$.
Hence
\begin{equation}\label{dsup:eq:allmultcohom}
 H^r(\Z^d,K^\times)
 \simeq H^r(\Z^d,M)\times\mm_F^{\binom dr}.
\end{equation}
This separates the fully computed principal-unit obstruction from the
ordinary finite-coordinate monomial-module problem. We do not claim that
$H^r(\Z^d,M)$ has a universal smaller formula for arbitrary shears and
characters.

\section{Rational dilations: only scalar obstructions survive}\label{dsup:sec:dilations}

Assume in this section that $\Gamma$ is divisible. Let
$q_1,\ldots,q_d\in\Q_{>0}$, with at least one $q_j\ne1$, and use
$\sigma_j=S_{q_j}$. Since a nonzero rational scalar acts injectively on
an ordered divisible group, the joint fixed exponent subgroup is $\{0\}$.
Therefore
\begin{equation}\label{dsup:eq:dilationfixed}
 F=\kk,\qquad P(x)=[x]_0,\qquad\mm_F=0.
\end{equation}

\begin{theorem}[Dilation equations and cohomology]\label{dsup:thm:dilationcohom}
For this family,
\begin{align}
 H^r(\Z^d,K_{\mathrm{add}})&\simeq\kk^{\binom dr},\label{dsup:eq:diladd}\\
 H^r(\Z^d,K^\times)&\simeq(\kk^\times)^{\binom dr}
 \quad(0\le r\le d).\label{dsup:eq:dilmult}
\end{align}
Both vanish for $r>d$. A compatible additive system is solvable exactly
when every right-hand side has constant coefficient zero. A compatible
multiplicative system is solvable exactly when every right-hand side
has leading coefficient one.
\end{theorem}
\begin{proof}
The additive formula and its solvability criterion follow from
\eqref{dsup:eq:dilationfixed}. In multiplicative coordinates, the monomial
module is now the direct product of $\Gamma$ acted on by rational
multiplication and $\kk^\times$ acted on trivially. On $\Gamma$, at least
one $q_j-1$ is an automorphism. Its inverse times the corresponding
exterior contraction contracts the entire Koszul complex. Thus all
cohomology of this $\Gamma$-module, including invariants, is zero.
The infinitesimal module is also acyclic by Corollary~\ref{dsup:cor:unitcohom}.
The trivial module $\kk^\times$ has zero Koszul differential, giving
\eqref{dsup:eq:dilmult}.

For an explicit check of the multiplicative criterion, choose $j_0$ with
$q_{j_0}\ne1$ and set $g=\gamma_{j_0}/(q_{j_0}-1)$. Compatibility gives
\[
 (q_{j_0}-1)\gamma_j=(q_j-1)\gamma_{j_0},
\]
so $\gamma_j=(q_j-1)g$ for every $j$. The character equations are exactly
$c_j=1$, and the logarithmic obstruction vanishes because $\mm_F=0$.
This proves both necessity and sufficiency without an appeal to the
abstract quotient description.
\end{proof}

For one $q\ne1$, these statements read
\begin{align}
 S_qy-y=f&\quad\Longleftrightarrow\quad [f]_0=0
       &&\text{for existence of }y\in K,\label{dsup:eq:scalaradd}\\
 S_qy/y=f&\quad\Longleftrightarrow\quad \lc(f)=1
       &&\text{for existence of }y\in K^\times.\label{dsup:eq:scalarmult}
\end{align}
The solutions are unique modulo addition of $\kk$ or multiplication by
$\kk^\times$, respectively. These are existence equivalences, not claims
that either equation has a solution for every right-hand side.

\begin{example}[Two opposite support directions]\label{dsup:ex:opposite}
For $S_2$ one has
\[
 (S_2-I)\left(-\sum_{n\ge0}t^{2^n}\right)=t,
 \qquad
 (S_2-I)\left(\sum_{n\ge1}t^{-2^{-n}}\right)=t^{-1}.
\]
The second solution has negative exponents increasing to zero. It is a
perfectly valid Hahn series, but it is not a Laurent series with a fixed
bound on exponent denominators. There is no solution to
$(S_2-I)y=1$.
\end{example}

\begin{example}[A multiplicative Mahler product]\label{dsup:ex:product}
The family defining
\[
 y=\prod_{n\ge0}(1-t^{2^n})^{-1}
\]
is interpreted either coefficientwise or by its strongly summable
logarithm. It satisfies $S_2y/y=1-t$: doubling drops the factor with
$n=0$. On the other hand, $S_2y/y=2$ has no nonzero solution. The
obstruction is its leading coefficient, not its valuation. These
particular rank-one identities are illustrative Mahler calculations,
not separate claims of novelty.
\end{example}

\section{Doubling defines the entire coefficient structure}\label{dsup:sec:definability}

For this section assume only that $\Gamma\ne0$ is $2$-divisible. (The
source abbreviated $S_2$ to $D$ in this section; the report writes $S_2$
throughout, see Section~\ref{dsup:sec:conventions}.) We work in the
one-sorted language
\[
 L_{S_2}=\{0,1,+,-,\cdot,S_2\}.
\]
Division by a nonzero element is a definable operation in this language.
Every auxiliary symbol introduced below abbreviates an explicit
first-order formula. Neither the coefficient field nor the value group
is named in $L_{S_2}$.

\begin{theorem}[Recovery from a single dilation]\label{dsup:thm:definability}
In $(K,S_2)$ the following data are parameter-free definable: the embedded
coefficient field, the constant-coefficient map, the canonical monomial
group $t^\Gamma$, the coefficient function on $K\times t^\Gamma$, the
support relation, the order of the monomial exponents, the valuation
ring and its maximal ideal, the leading monomial and coefficient, and
truncation at any monomial exponent.
\end{theorem}

The proof is constructive. We give the definitions and check each one.

\subsection{Constants and the constant-coefficient map}

By Corollary~\ref{dsup:cor:fixed}, since $2g=g$ only for $g=0$,
\begin{equation}\label{dsup:eq:defconstants}
 C(c):\quad S_2c=c
\end{equation}
defines precisely $\kk$. Theorem~\ref{dsup:thm:resolvent} gives
\begin{equation}\label{dsup:eq:constantformula}
 B(x,c):\quad S_2c=c\ \land\ \exists y\;(x-c=S_2y-y).
\end{equation}
For every $x$ there is exactly one $c$ satisfying $B(x,c)$, namely $[x]_0$.
Write $B(x)$ for this definable function. Only the inverse of doubling
on $\Gamma$ is needed in that application of the resolvent theorem.

\subsection{The equation that recognizes monomials}

\begin{lemma}[Monomial recognition]\label{dsup:lem:monomialrecognition}
For $x\in K$,
\begin{equation}\label{dsup:eq:monomialformula}
 x\in t^\Gamma\quad\Longleftrightarrow\quad x\ne0\ \land\ S_2x=x^2.
\end{equation}
\end{lemma}
\begin{proof}
Every $t^g$ satisfies $S_2(t^g)=t^{2g}=(t^g)^2$.
Conversely, let $x\ne0$ satisfy the equation. Write
$x=ct^g(1+u)$, where $c\in\kk^\times$ and $u\in\mm$. Comparison of
leading coefficients gives $c=c^2$, so $c=1$. The remaining equation is
$S_2(1+u)=(1+u)^2$. If $u\ne0$, put $\eps=v(u)>0$.
The left-hand side minus $1$ has valuation $2\eps$, whereas
$2u+u^2$ has valuation $\eps$, since the characteristic is zero and
$v(u^2)=2\eps>\eps$. This is impossible. Thus $u=0$ and $x=t^g$.
\end{proof}

Let $M(m)$ denote the formula on the right of
\eqref{dsup:eq:monomialformula}. An important feature is that it selects
coefficient \emph{one}, not arbitrary nonzero scalar multiples of
monomials. It is also entirely algebraic; no positivity condition on
$x$ occurs.

\subsection{All coefficients and the support relation}

For $M(m)$ define
\begin{equation}\label{dsup:eq:coefficientformula}
 \coef(x,m)=B(x/m).
\end{equation}
If $m=t^g$, this equals $x_g$. Hence
\begin{equation}\label{dsup:eq:supportformula}
 \Supp(m,x):\quad M(m)\ \land\ \coef(x,m)\ne0
\end{equation}
is the exact support relation expressed on monomials. In particular,
\begin{equation}\label{dsup:eq:coefficientsseparate}
 x=y\quad\Longleftrightarrow\quad
 \forall m\bigl(M(m)\Rightarrow\coef(x,m)=\coef(y,m)\bigr).
\end{equation}
This step recovers the full series, not merely its leading term.

\subsection{Recovering the orientation of the valuation}

Define
\begin{equation}\label{dsup:eq:infmon}
 J(m):\quad M(m)\ \land\ m\ne1\ \land\ B((1-m)^{-1})=1.
\end{equation}
For $m=t^g$ with $g>0$, the expansion
$(1-m)^{-1}=\sum_{n\ge0}m^n$ has constant coefficient one.
For $g<0$, the correct expansion is
\[
 (1-m)^{-1}=-\sum_{n\ge1}m^{-n},
\]
whose constant coefficient is zero. Both supports are well ordered in
their indicated directions. Thus
\begin{equation}\label{dsup:eq:infmeaning}
 J(t^g)\quad\Longleftrightarrow\quad g>0.
\end{equation}
Consequently, for monomials $m,n$,
\begin{equation}\label{dsup:eq:monomialorder}
 v(m)<v(n)\quad\Longleftrightarrow\quad J(n/m).
\end{equation}
Multiplication of monomials corresponds to addition in $\Gamma$, so
this defines the entire ordered value group on the definable set $M$.
For real coefficients its order is the reverse of the usual order of
positive monomials. For complex coefficients no field order is required.

\subsection{The valuation ring, leading term, and truncation}

The following formulas now define $\OO$ and $\mm$:
\begin{align}
 x\in\OO&\quad\Longleftrightarrow\quad
  \forall m\bigl(\Supp(m,x)\Rightarrow(m=1\ \lor\ J(m))\bigr),\label{dsup:eq:defO}\\
 x\in\mm&\quad\Longleftrightarrow\quad
  \forall m\bigl(\Supp(m,x)\Rightarrow J(m)\bigr).\label{dsup:eq:defm}
\end{align}
They include zero by vacuous quantification, as they should. For $x\ne0$,
the leading monomial is the unique $\ell$ such that
\begin{equation}\label{dsup:eq:leadingformula}
 \Supp(\ell,x)\ \land\
 \forall m\bigl(\Supp(m,x)\Rightarrow(m=\ell\ \lor\ J(m/\ell))\bigr).
\end{equation}
Its coefficient is $\coef(x,\ell)$. Existence follows from the
well-ordering of the support.

For $M(m)$, the truncation $T(x,m)$ removes all exponents at least $v(m)$:
\[
 T(x,t^g)=\sum_{h<g}x_ht^h.
\]
Its graph is given by
\begin{equation}\label{dsup:eq:truncformula}
 \begin{split}
 y=T(x,m)\quad\Longleftrightarrow\quad M(m)\ \land\
 \forall n\bigl(M(n)\Rightarrow{}&
 [(J(m/n)\Rightarrow\coef(y,n)=\coef(x,n))\ \land\ {}\\[-2pt]
 & (\neg J(m/n)\Rightarrow\coef(y,n)=0)]\bigr).
 \end{split}
\end{equation}
The retained support is a subset of a well-ordered set, so the required
$y$ exists. Equation~\eqref{dsup:eq:coefficientsseparate} gives uniqueness.
This completes the proof of Theorem~\ref{dsup:thm:definability}.

\begin{remark}[The role of fullness]
The proof uses that the orbit sums solving $S_2y-y=f-[f]_0$ belong to the
field. An arbitrary subfield of $\kk((t^\Gamma))$ that is stable under
$S_2$ need not have this property. Likewise, all Boolean subseries used
below must exist. Thus replacing the full Hahn field by a rational
function field, a finite-support ring, or an arbitrary union of Laurent
fields requires new hypotheses. Definability in the full field does not
supply an algorithm for arbitrary presentations of its elements.
\end{remark}

\section{The full centralizer of doubling}\label{dsup:sec:centralizer}

The defining formulas have a strong algebraic consequence. A field
automorphism commuting with doubling must preserve all the recovered
normal-form data, even when nothing of that kind is assumed beforehand.

\begin{theorem}[Unrestricted centralizer]\label{dsup:thm:centralizer}
Let $\Gamma\ne0$ be $2$-divisible. Every field automorphism
$\alpha\in\Aut(K)$ satisfying $\alpha S_2=S_2\alpha$ has a unique form
\begin{equation}\label{dsup:eq:centralizerform}
 \alpha\left(\sum_g a_gt^g\right)
       =\sum_g\rho(a_g)t^{\theta g},
 \quad \rho\in\Aut(\kk),\quad
       \theta\in\Aut_{\mathrm{ord}}(\Gamma).
\end{equation}
Conversely, every map \eqref{dsup:eq:centralizerform} commutes with $S_2$.
Therefore
\begin{equation}\label{dsup:eq:centralizer}
 \Cen_{\Aut(K)}(S_2)
     \simeq\Aut(\kk)\times\Aut_{\mathrm{ord}}(\Gamma).
\end{equation}
Every member of this centralizer is strong and preserves the valuation
ring and all truncations.
\end{theorem}
\begin{proof}
Since $\alpha$ is an automorphism of the $L_{S_2}$-structure, it preserves
every parameter-free definable set and function in
Theorem~\ref{dsup:thm:definability}. In particular its restriction to the
coefficient field is an automorphism $\rho$ of $\kk$.
It permutes the canonical monomials, so write
$\alpha(t^g)=t^{\theta g}$. Multiplicativity makes $\theta$ additive;
bijectivity makes it a group automorphism. Preservation of $J$ gives
$g>0\Longleftrightarrow\theta g>0$, hence $\theta$ is increasing.

For every $x\in K$ and $g\in\Gamma$, preservation of the coefficient
function gives
\[
 \coef(\alpha x,t^{\theta g})
 =\coef(\alpha x,\alpha(t^g))
 =\alpha(\coef(x,t^g))=\rho(x_g).
\]
The map $\theta$ is onto, so this computes every coefficient of $\alpha x$.
It proves \eqref{dsup:eq:centralizerform} by uniqueness of Hahn expansions.
There is no prior assumption that $\alpha$ commutes with infinite sums.
The formula itself now proves that it does: $\theta$ preserves
well-ordered supports, while $\rho$ preserves finite sums and nonzero
coefficients.

Conversely, \eqref{dsup:eq:centralizerform} defines a strong field automorphism.
Additivity of $\theta$ gives $\theta(2g)=2\theta(g)$, so it commutes with
$S_2$. The two choices are unique on coefficients and monomials. Their
composition is componentwise, since coefficient and exponent operations
commute, which gives the direct product in \eqref{dsup:eq:centralizer}.
\end{proof}

\begin{corollary}[No hidden centralizing motions]\label{dsup:cor:nohidden}
A field automorphism commuting with $S_2$ and fixing every coefficient
and canonical monomial is the identity. In particular a centralizing
automorphism which fixes every leading term is the identity.
\end{corollary}
\begin{proof}
For the first assertion use \eqref{dsup:eq:centralizerform} with
$\rho=\id$ and $\theta=\id$. For the second, the formula implies that
fixing all leading terms fixes the coefficient and exponent components,
so the first assertion applies.
\end{proof}

A character twist $t^g\mapsto\chi(g)t^g$ illustrates the force of the
condition: it commutes with doubling only if
$\chi(2g)=\chi(g)$, but the character law says
$\chi(2g)=\chi(g)^2$. Thus $\chi(g)=1$ for every $g$.
The theorem also excludes centralizing automorphisms which are not
monomial twists and might initially have been nonvalued or nonstrong.

When $\kk=\R$, $\Aut(\kk)$ is trivial. Indeed an automorphism of $\R$
preserves positivity because positive elements are exactly nonzero
squares, fixes $\Q$, and therefore fixes each real number by its rational
cut. When $\kk=\C$, arbitrary field automorphisms of $\C$ remain; only
those equal to the identity or complex conjugation preserve $\R$ setwise,
by the preceding observation and $\C=\R(i)$.

\section{Arithmetic interpreted by one automorphism}\label{dsup:sec:logic}

This section identifies the logical cost of recovering normal-form
information. The implication from support to arithmetic is related to
Camacho's theorem \cite{Camacho}; the formulas below show directly how it
is implemented in the smaller, single-dilation signature. The elementary
coding is included so the dependence on full Hahn supports is explicit.

\subsection{A uniform Boolean-series interpretation}

Choose a parameter $m$ satisfying $J(m)$; such a monomial exists since
$\Gamma\ne0$. Set $b_m=(1-m)^{-1}$. Its support is exactly
$\{m^n:n\in\N\}$. Define
\[
 N_m=\{u:M(u)\text{ and }\coef(b_m,u)\ne0\}.
\]
Field multiplication on $N_m$ represents addition of natural numbers:
$m^r m^s=m^{r+s}$. The represented zero and one are $1$ and $m$.
Define the Boolean-series set
\begin{equation}\label{dsup:eq:Booleanseries}
 \begin{split}
 \mathcal B_m=\{x\in K:\ \forall u\,[M(u)\Rightarrow{}&
 (\coef(x,u)=0\ \lor\ \coef(x,u)=1)]\ \land{}\\[-2pt]
 &\forall u\,[\Supp(u,x)\Rightarrow u\in N_m]\}.
 \end{split}
\end{equation}
Each subset $A\subseteq\N$ has exactly one representative
$x_A=\sum_{n\in A}m^n$ in $\mathcal B_m$. Conversely every member of
$\mathcal B_m$ has this form. Membership is the definable relation
$\coef(x_A,m^n)=1$. Thus the two definable sorts $N_m,\mathcal B_m$,
with these operations and relation, are exactly
$(\N,\mathcal P(\N);+,\in)$. No quotient by equal supports is needed
because Boolean coefficients make the coding unique.

\begin{theorem}[Undecidability in the unexpanded signature]\label{dsup:thm:undecidable}
For every nonzero $2$-divisible $\Gamma$ and characteristic-zero field
$\kk$, the complete first-order theory of $(\kk((t^\Gamma)),S_2)$ is
undecidable. It uniformly interprets full monadic second-order arithmetic
with one suitable monomial parameter. The undecidability assertion does
not require adding a constant symbol for that parameter.
\end{theorem}
\begin{proof}
The preceding definitions give an effective translation of every
sentence $\varphi$ of $(\N,\mathcal P(\N);+,\in)$ to a formula
$\varphi^m$ with the parameter $m$. Every $m$ satisfying $J$ yields the
same standard interpreted structure. Therefore
\begin{equation}\label{dsup:eq:parameterremoval}
 (\N,\mathcal P(\N);+,\in)\models\varphi
 \quad\Longleftrightarrow\quad
 (K,S_2)\models\exists m\,(J(m)\land\varphi^m).
\end{equation}
For completeness, multiplication on $\N$ is definable in the interpreted
structure. The relation $a\mid b$ says that every subset containing zero
and closed under addition by $a$ also contains $b$. This is exactly
membership in $a\N$, because that set is the least such subset.
For $a>0$, the least positive common multiple of $a$ and $a+1$ is
$a(a+1)$: consecutive positive integers are coprime. Leastness is
expressible with the order defined from addition. For $a=0$ define this
product to be zero. Its difference from $a$ defines $a^2$. Finally
\[
 (a+b)^2=a^2+b^2+2ab
\]
defines $ab$ uniquely using addition and squares. Hence ordinary true
arithmetic is interpreted and its undecidability rules out a decision
algorithm for the theory of $(K,S_2)$ via
\eqref{dsup:eq:parameterremoval}. The explicit existential quantification is
important: interpretation with an arbitrary named parameter alone would
not justify this particular reduct claim.
\end{proof}

The interpreted second sort is genuinely all of $\mathcal P(\N)$, not
just computable sets or eventually periodic sets. This is why full Hahn
series matter. The theorem concerns the complete theory of the specified
structure; it does not assert a failure of decidability for every smaller
language associated with a valued field.

\subsection{A direct coefficient witness to \texorpdfstring{$\mathrm{TP}_2$}{TP2}}

The following argument is independent of the arithmetic reduction and
uses arbitrary coefficients rather than only zero and one.

\begin{theorem}[Coefficient independence]\label{dsup:thm:tp2}
The formula
\[
 \varphi(x;u,c):\quad M(u)\ \land\ C(c)\ \land\ \coef(x,u)=c
\]
has the tree property of the second kind in the complete theory of
$(K,S_2)$. The support relation also has the independence property and
the strict order property.
\end{theorem}
\begin{proof}
Fix $m$ with $J(m)$ and use parameters
$(u_i,c_j)=(m^i,j)$ for $i,j\in\N$, where $j$ is its characteristic-zero
image in $\kk$. For each fixed row $i$, the formulas
$\varphi(x;m^i,j)$ for different $j$ are pairwise inconsistent: a
coefficient has only one value. For every function
$\eta:\N\to\N$, the Hahn series
\[
 x_\eta=\sum_{i\ge0}\eta(i)m^i
\]
exists and satisfies all formulas
$\varphi(x;m^i,\eta(i))$. Its support is a subset of the well-ordered
set $\{i\,v(m):i\in\N\}$, regardless of the sizes of its coefficients.
This array is exactly a $\mathrm{TP}_2$ witness, with two-inconsistent
rows and every choice of one entry in each row realized.

The Boolean series $x_A$ above shatter the infinite set
$\{m^i:i\in\N\}$ under the support relation, giving the independence
property. Finally the supports of
$1+m+\cdots+m^n$ form a strictly increasing chain of definable sets of
monomials. This witnesses the strict order property.
\end{proof}

\begin{remark}[No contradiction with tame valued-field reducts]
Adding $S_2$ is a genuine expansion. The formulas use $S_2$ to recover
coefficient extraction throughout the field, including its purely
infinite part. The resulting interpreted arithmetic does not follow
merely from naming a valuation ring. Likewise, this coefficient-fixing
structure is not being identified with the existentially closed residue
difference structures occurring in other theories of valued difference
fields \cite{Pal}.
\end{remark}

\section{Every nonidentity rational dilation recovers the same data}\label{dsup:sec:alldilations}

Assume again that $\Gamma\ne0$ is divisible. For
$q=a/b\in\Q_{>0}\setminus\{1\}$, with positive integers $a,b$, let
$n=|a-b|>0$ and put
\[
 E_{a,b}(x):\quad x\ne0\ \land\ S_q(x)^b=x^a.
\]
The next result both extends the recovery theorem and distinguishes
interdefinability from conjugacy.

\begin{lemma}[Power-equation recognition]\label{dsup:lem:rationalrecognition}
In $(K,S_q)$,
\[
 \{x:E_{a,b}(x)\}=\mu_n(\kk)\,t^\Gamma,
 \qquad
 t^\Gamma=\{y^n:E_{a,b}(y)\},
\]
where $\mu_n(\kk)=\{c\in\kk^\times:c^n=1\}$.
\end{lemma}
\begin{proof}
Write $x=ct^g\exp(h)$. The exponent equation is automatic since
$bq=a$. The coefficient equation is $c^b=c^a$, equivalently $c^n=1$.
The principal-unit equation is $bS_qh=ah$, or $S_qh=qh$.
The resonance set for the scalar $\lambda=q$ is empty: the only fixed
exponent is zero and its character value is $1\ne q$. Therefore
Theorem~\ref{dsup:thm:resolvent} implies $h=0$. This proves the first equality.
The $n$th powers of these elements are precisely $t^{n\Gamma}=t^\Gamma$,
because $\Gamma$ is divisible and every $c\in\mu_n(\kk)$ has $c^n=1$.
\end{proof}

\begin{theorem}[Interdefinable dilations]\label{dsup:thm:rationaldef}
For each nonidentity $q\in\Q_{>0}$, the structure $(K,S_q)$ parameter-free
defines all the normal-form data of Theorem~\ref{dsup:thm:definability} and
every rational dilation $S_r$. In particular,
\[
 \Cen_{\Aut(K)}(S_q)
 =\{\text{untwisted coefficient--ordered-exponent lifts}\}
 \simeq\Aut(\kk)\times\Aut_{\mathrm{ord}}(\Gamma).
\]
These centralizers are the same subgroup of $\Aut(K)$ for all $q\ne1$.
The arithmetic and $\mathrm{TP}_2$ conclusions hold for every such $S_q$.
\end{theorem}
\begin{proof}
The fixed field of $S_q$ is $\kk$, and
$x-c\in\im(S_q-I)$ defines $c=[x]_0$, by
\eqref{dsup:eq:scalaradd}. Lemma~\ref{dsup:lem:rationalrecognition} defines the
canonical monomial set. Once these two pieces are known, every subsequent
formula in Section~\ref{dsup:sec:definability} is unchanged.

For a positive rational $r=c/d$, the monomial power $m^r$ is the unique
monomial $z$ satisfying $z^d=m^c$; existence uses divisibility of
$\Gamma$ and uniqueness uses its torsion-freeness. The graph of $S_r$ is
then defined by
\[
 y=S_rx\quad\Longleftrightarrow\quad
 \forall m\bigl(M(m)\Rightarrow
       \coef(y,m^r)=\coef(x,m)\bigr).
\]
Every exponent is an $r$-multiple, so this determines all coefficients.
Existence follows from the increasing exponent automorphism $g\mapsto rg$.
In particular $S_2$ is definable from $S_q$, and conversely $S_q$ is
definable from $S_2$. The centralizer statement follows either from these
mutual definitions or by repeating the proof of
Theorem~\ref{dsup:thm:centralizer}. The logical consequences transfer through
the explicit definitions.
\end{proof}

\begin{theorem}[Distinct rational dilations are not conjugate]\label{dsup:thm:nonconjugate}
If $r,s\in\Q_{>0}$ are distinct, then $S_r$ and $S_s$ are not conjugate
by any field automorphism of $K$. Moreover,
\[
 \Nor_{\Aut(K)}(\langle S_2\rangle)
   =\Cen_{\Aut(K)}(S_2).
\]
\end{theorem}
\begin{proof}
The identity dilation is not conjugate to a nonidentity map, so assume
$r\ne1$. Write $r=a/b$ and $n=|a-b|$. The equation
\[
 S_r(x)^b=x^a,\qquad x\ne0,
\]
has infinitely many solutions, since every $t^g$ is a solution and
$\Gamma\ne0$ is infinite. For $S_s$ in place of $S_r$, the valuation
comparison gives $(bs-a)v(x)=0$, so $v(x)=0$ because $s\ne r$.
Writing $x=c\exp(h)$ reduces the equation to $c^n=1$ and $S_sh=rh$.
There is no nonzero solution of the latter scalar eigen-equation:
if $s\ne1$, its only possible fixed exponent is zero, with eigenvalue
$1\ne r$; if $s=1$, it is simply $(1-r)h=0$.
Thus there are at most $n$ nonzero solutions for $S_s$.
Conjugation would carry the first solution set bijectively to the second,
a contradiction. Equivalently, a single first-order sentence requiring
more than $n$ distinct solutions distinguishes the two pointed fields.

The group $\langle S_2\rangle$ is infinite cyclic. An automorphism
normalizing it must send its generator to $S_2$ or $S_2^{-1}=S_{1/2}$.
The second alternative is excluded by nonconjugacy, leaving exactly the
centralizer.
\end{proof}

Interdefinability and nonconjugacy are compatible. Interdefinability
recovers each particular dilation by a formula from another; conjugacy
would have to carry the distinguished operator itself to the other one.
The power equations distinguish the latter requirement.

\section{Arbitrary-rank examples and sharp distinctions}\label{dsup:sec:examples}

\subsection{A shear with a large fixed Hahn field}

Take $\Gamma=\Q^2$ in lexicographic order and
\[
 \theta(a,b)=(a,b+a),\qquad\chi=1.
\]
This is an increasing group automorphism, with inverse
$(a,b)\mapsto(a,b-a)$. Its increasing and decreasing chambers are $a>0$
and $a<0$, respectively. The fixed exponent subgroup is
$\{0\}\times\Q$, so
\[
 K^\sigma=\kk((t^{\{0\}\times\Q})).
\]
The entire series supported on that subgroup is the additive obstruction,
not just the scalar constant coefficient. For example,
\begin{equation}\label{dsup:eq:shear}
 (\sigma-I)\left(-\sum_{n\ge0}t^{(1,n)}\right)=t^{(1,0)},
 \qquad
 -\sum_{n\ge0}t^{(1,n)}
   =-\frac{t^{(1,0)}}{1-t^{(0,1)}}.
\end{equation}
The inverse is exact although the exponents $(1,n)$ never exceed $(2,0)$.
The errors in finite partial sums therefore do not eventually have
valuation greater than $(2,0)$. Equation~\eqref{dsup:eq:shear} is a concrete
reason to retain strong support arguments rather than impose a rank-one
analytic convergence intuition.

The additive and multiplicative obstructions are genuinely different.
A nonzero fixed monomial $t^{(0,b)}$ cannot be in the additive image of
$\sigma-I$. Nevertheless it is a multiplicative coboundary:
\[
 \frac{\sigma(t^{(b,0)})}{t^{(b,0)}}=t^{(0,b)}.
\]
Being fixed does not itself prevent multiplicative exactness.

\subsection{A phase character and an integral-exponent obstruction}

Take $K=\C((t^\Q))$, $\theta=\id$, and
$\chi(g)=e^{2\pi i g}$. Then
\[
 \sigma\left(\sum_g a_gt^g\right)
       =\sum_g a_ge^{2\pi i g}t^g,
 \qquad K^\sigma=\C((t^\Z)).
\]
Here every exponent is fixed by $\theta$, so solving
$(\sigma-I)y=f$ means deleting no moving orbits at all. It is possible
exactly when every integral-exponent coefficient of $f$ vanishes; the
other coefficients are divided by $e^{2\pi i g}-1$.

These denominators may be arbitrarily small as ordinary complex numbers.
They do not obstruct strong Hahn summability: no support points are
introduced and no infinite coefficientwise sums occur. Thus analytic
small-denominator estimates and Hahn support obstructions are different
issues.

For a multiplicative right-hand side $f=ct^\gamma\exp(h)$, the criterion
is
\[
 \gamma=0,\qquad c\text{ is a root of unity},\qquad
 \restr h{\Z}=0.
\]
Indeed $(\theta-I)g=0$, and $\chi(\Q)$ is exactly the group of complex
roots of unity. In the last condition $h\in\mm$, so only positive
integral exponents can occur. The ambient field is algebraically closed by the standard full-Hahn
closedness theorem for algebraically closed characteristic-zero coefficients
and a divisible value group; the repository also documents this theorem
in its Hahn layer \cite{Repo}. The fixed field is not algebraically
closed: a square root of $t$ would have valuation $1/2$, which is absent
from $\Z$. Thus the fixed field in the cohomology calculation can lose
algebraic closedness.

\subsection{Compatibility is not the same as exactness}

For any commuting family, compatibility only says that a tuple is a
cocycle. The joint-resonant coefficients need not vanish. A simple
additive example is a nonzero tuple of constants: it is automatically
compatible, but no nonzero entry can occur in an exact tuple for a
nonidentity rational dilation family.

A similarly simple multiplicative example is a tuple
$(c_1,\ldots,c_d)\in(\kk^\times)^d$. It is compatible because the
coefficient action is trivial. For a rational-dilation family with at
least one nonidentity member, it is exact precisely when every $c_j=1$.
These examples realize every scalar obstruction appearing in
Theorem~\ref{dsup:thm:dilationcohom}.

\section{Transfer to surreal and surcomplex normal forms}\label{dsup:sec:surreal}

\subsection{The normal-form dictionary}

The standard Conway normal form writes a surreal as
\[
 x=\sum_{\beta<\alpha}r_\beta\omega^{a_\beta},
 \qquad r_\beta\in\R,
 \quad (a_\beta)_{\beta<\alpha}\text{ strictly decreasing},
\]
where $\alpha$ is an ordinal. The normal-form theorem, and its use as a
Hahn-field representation, is described in \cite{KKS}. With
\begin{equation}\label{dsup:eq:surrealdictionary}
 t^g=\omega^{-g},
\end{equation}
the support is well ordered in increasing $g$-notation. Combining the
supports of real and imaginary parts gives the corresponding
complex-coefficient normal form for $\SC=\No(i)$.

Thus, for $q\in\Q_{>0}$, define the class field automorphisms
\begin{equation}\label{dsup:eq:surrealdilation}
 S_q\left(\sum_{\beta<\alpha}c_\beta\omega^{a_\beta}\right)
       =\sum_{\beta<\alpha}c_\beta\omega^{qa_\beta},
 \quad c_\beta\in\R\text{ or }\C.
\end{equation}
These are precisely the usual coefficient-fixing monomial dilations;
their construction is not new. The new assertions concern what their
difference equations and their centralizers determine.

\subsection{Set-sized localization of every proof}

A set of normal forms uses a set of exponents. For a finite family of
class exponent automorphisms, close that set under their positive and
negative integer iterates and under finite group operations. This
closure is still a set. When division by rational scalars is needed,
take its rational span. The resulting ordered subgroup supplies a
set-sized Hahn field containing all input supports and all the orbit
supports used in the additive constructions. Adding any exponent $g$
needed for a multiplicative monomial equation still gives a set-sized
subgroup. Neumann sums use finite sums of already included exponents.

More quantitatively, if the original union of supports has cardinal
$\kappa$, these support-enlargement operations require at most
$\max(\kappa,\aleph_0)$ exponents, apart from a finite number of explicitly
adjoined monomial-solution exponents. This does not bound their surreal
birthdays and does not assert that a chosen birthday fragment is closed
under all these operations.

The ordinary cohomology groups in Theorems~\ref{dsup:thm:koszul} and
\ref{dsup:thm:dilationcohom} remain statements about these set-sized fields.
For the full proper-class fields, we state the corresponding
\emph{elementwise equation and formula assertions}. They involve
set-supported series and finite tuples, not a purported set of all
class cochains.

\begin{theorem}[Surreal and surcomplex consequences]\label{dsup:thm:surreal}
Let $\mathbb F$ be $\No$ or $\SC$, with coefficient field $\kk=\R$ or
$\C$, respectively. Then the following assertions hold.
\begin{enumerate}[leftmargin=*,itemsep=4pt]
\item For $q\in\Q_{>0}\setminus\{1\}$,
$S_qy-y=f$ is solvable in $\mathbb F$ exactly when the coefficient of
$\omega^0$ in $f$ is zero. The solution is unique modulo $\kk$.
\item For $f\in\mathbb F^\times$, the equation $S_qy/y=f$ is solvable
exactly when the leading normal-form coefficient of $f$ is one. The
solution is unique modulo $\kk^\times$.
\item The simultaneous criteria of Theorem~\ref{dsup:thm:dilationcohom} hold
for every finite family of positive rational dilations with at least
one nonidentity member.
\item In the field expanded by $S_q$, the formulas in
Sections~\ref{dsup:sec:definability} and \ref{dsup:sec:alldilations} define the
ordinary coefficient field, the Conway monomial class, every normal-form
coefficient, the natural valuation ring, and all monomial truncations.
\item If a class field automorphism $\alpha$ commutes with $S_q$, then
there are unique $\rho\in\Aut(\kk)$ and an ordered additive class
automorphism $\theta$ of $(\No,+,<)$ such that
\begin{equation}\label{dsup:eq:surrealcentralizer}
 \alpha\left(\sum_{\beta<\alpha_0}c_\beta\omega^{a_\beta}\right)
      =\sum_{\beta<\alpha_0}\rho(c_\beta)\omega^{\theta(a_\beta)}.
\end{equation}
Conversely every such class map commutes with $S_q$. For $\No$, necessarily
$\rho=\id_\R$.
\end{enumerate}
\end{theorem}
\begin{proof}
For the equation assertions, put the given supports into a set-sized
Hahn workspace as just described and use the explicit solutions already
proved. Necessity can also be checked directly on full normal forms;
allowing larger supports does not remove a resonant coefficient
obstruction.

The defining formulas hold for all surreal or surcomplex arguments,
because their proofs use only the unique normal form of the element
being tested and explicit set-supported witnesses. For the centralizer,
repeat the preservation-of-definable-coefficients argument from
Theorem~\ref{dsup:thm:centralizer}. The normal-form monomial group is identified
with $(\No,+)$ by \eqref{dsup:eq:surrealdictionary}, and the recovered valuation
orientation gives its order. Additivity of $\theta$ makes the switch from
$g$ to $a=-g$ in \eqref{dsup:eq:surrealcentralizer} consistent. Every class
image of a set support is again a set in a class theory with replacement
for the class maps under consideration. Uniqueness of normal forms
therefore gives the formula for each element. The converse is the
standard strong monomial-lift calculation. Finally a field automorphism
of $\R$ is the identity, as proved after
Corollary~\ref{dsup:cor:nohidden}.
\end{proof}

\begin{remark}[Foundational reading]
These full-field statements can be read in NBG-style class mathematics,
with each assertion about a class automorphism made for a specified
class map. The notation for a centralizer of class automorphisms is
structural shorthand, not an assertion that all such automorphisms form
a set or even an ordinary class of set objects. All ordinary group
cohomology and complete first-order theory claims in this article have
set-sized versions and proofs. The same finite interpretation formulas
also encode ordinary arithmetic and its power-set sort in the full
surreal expansions; no global satisfaction class for arbitrary class
structures is needed for those individual formula identities.
\end{remark}

\subsection{What is recovered, and what is not}\label{dsup:sub:recovered}

On $\SC$, the theorem recovers the \emph{complex} coefficient field and
the canonical monomials. It does not claim that the distinguished real
axis $\No$ is thereby definable. The centralizer still allows arbitrary
coefficient-field automorphisms of $\C$; the ones preserving $\R$ are
only the identity and conjugation. Thus preserving the recovered
valuation structure is not the same as preserving a chosen real form.

Nor does $S_q$ here mean the surreal exponential, the omega-map itself,
a Berarducci--Mantova derivation, or an automorphism known to preserve
any of those additional structures. The centralizer theorem is a result
about an explicitly named \emph{field} automorphism. It does not settle
the exponential-automorphism image problems discussed in \cite{KKS,RepoAut}.
Finally the full-surreal fine topology is not the summation mechanism
used in the proofs. The repository distinguishes its small-net behavior
from strong Hahn summation \cite{RepoDocs}; the present constructions stay
on the latter side of that distinction.

\section{A formalization route and precise limitations}\label{dsup:sec:formalization}

\subsection{A dependency-oriented route}

A Lean development can be organized without attempting to formalize the
model theory and the proper-class interpretation at the same time as the
support algebra. The natural first target is a set-sized Hahn field,
using the repository's documented summability interface \cite{Repo}.

The first layer is order-theoretic: prove
Lemma~\ref{dsup:lem:orbit} as a summability certificate for the map
$(n,s)\mapsto\theta^ns$ on the increasing chamber, and separately prove
finite fibers. This should not be replaced by a convergence theorem.
The labeled version then permits the construction of a strong linear
$Q_\lambda$ from the two iterate families and the fixed-exponent quotient.

The next layer establishes the two inverse identities, normalization,
and the support/valuation bound. Once uniqueness of normalized solutions
is available, the commutation lemma can be proved abstractly and need
not repeat the coefficient calculation for every pair of automorphisms.
The Koszul contraction is then finite exterior algebra with commuting
linear maps. Its only Hahn-specific input is the previously verified
resolvent package.

For multiplicative equations, the existing infinitesimal Hahn
exponential/logarithm interface is the appropriate dependency. It is
important that the development type the group $U_1$ separately from the
whole multiplicative field, and keep the monomial-coordinate quotient
as an explicit residual obstruction.

The definability theorem can first be recorded as algebraic equivalences:
$S_2x=x^2$ iff $x$ is a monomial; existence of $y$ in $x-c=S_2y-y$ iff $c$ is
the constant coefficient; and the two geometric expansions deciding
$B((1-m)^{-1})$. A later model-theory layer can translate these
identities into formula syntax. The unrestricted centralizer theorem
then follows by preservation of these definable operations, rather than
by assuming a strong automorphism instance at the outset.

\subsection{Claims deliberately not made}\label{dsup:sub:notmade}

There is no Lean source accompanying this manuscript and no claim of
machine-checked infinite mathematics. The finite Python tests check signs,
endpoint identities, and exterior-algebra contractions only. There is no
uniform algorithm here for arbitrary Hahn names: deciding an exponent's
movement chamber or testing exact equality of coefficient values may
itself require unavailable computational information.

The coefficient action in \eqref{dsup:eq:sigma} is the identity. Allowing a
nontrivial difference automorphism of $\kk$ changes the fixed-exponent
recurrences; division by $\chi(g)-\lambda$ is no longer the right
operation. The finite family is commuting; no analogous contraction is
asserted for a nonabelian acting group. The multiplicative results are
abelian $K^\times$ results, not matrix gauge-classification theorems.
No support-growth, convergence, or solvability theorem for arbitrary
variable-coefficient Mahler systems is inferred from the scalar
resolvent.

The use of full Hahn fields is substantive. A restricted support class
must be checked for closure under increasing half-orbits, infinitesimal
logarithms, and all relevant subseries before any theorem is transferred.
Finally, characteristic zero is retained throughout the article. In
characteristic two the monomial-recognition equation degenerates when
coefficient squaring is trivial, so that particular formula cannot be
silently exported.

\subsection{Consolidated ledger of non-claims}\label{dsup:sec:nonclaims}

The source states its limitations where they arise, and each is kept in
place. The list collects them, with the place where each stands.
Items~1--28 are the source's; items~29--33 were added when the report
joined the collection.
\begin{enumerate}[label=\arabic*.,leftmargin=2.2em,itemsep=2pt]
\item Not refereed. No theorem is Lean-verified, no Lean source accompanies
the report, and nothing is represented as already checked by Lean; the
formalization route of Section~\ref{dsup:sec:formalization} is a route, not a
completed formalization (title page, Section~\ref{dsup:sec:intro},
Section~\ref{dsup:sub:notmade}, Appendix~\ref{dsup:app:provenance}; the
delivered README).
\item No named published open problem is claimed settled. In particular
nothing here determines the existence or image of nontrivial exponential
automorphisms of $\No$, and Camacho's monomial-definability question in the
smaller language without the extra automorphism is not answered (title
page, Section~\ref{dsup:sec:intro}, Section~\ref{dsup:sub:recovered},
Appendix~\ref{dsup:app:provenance}).
\item Priority is not certified. The main statements are proposed additions;
the literature review was focused (Hahn automorphisms, Mahler solutions,
valued difference fields, truncation definability), not exhaustive, and a
later-located antecedent would change the priority assessment, not the
proofs (title page, Section~\ref{dsup:sec:intro},
Appendix~\ref{dsup:app:provenance}).
\item The repository was not independently built or exhaustively audited:
only the catalogue, the root README and the surcomplex-automorphisms
report's guide were inspected, and that report's full article was not
audited (Section~\ref{dsup:sec:intro}, bibliography entry \cite{RepoAut}; the
delivered README).
\item Credited, not claimed: monomial and character lifts (Kaplan--Krapp--Serra
and the surcomplex-automorphisms report); the construction of the dilations
on $\No$ and $\No(i)$; Hahn field operations and the normal-form dictionary;
the support-to-arithmetic principle (Camacho); the principal-unit logarithm
and exponential and the Neumann lemma (Sections~\ref{dsup:sec:intro},
\ref{dsup:sec:multiplicative}, \ref{dsup:sec:logic}, \ref{dsup:sec:surreal},
Appendix~\ref{dsup:app:provenance}).
\item The rank-one Mahler identities of Examples~\ref{dsup:ex:opposite}
and~\ref{dsup:ex:product} are illustrative calculations, not separate
claims of novelty (Section~\ref{dsup:sec:dilations}).
\item The first-order, constant-scalar resolvent does not replace the
higher-order, variable-coefficient theory of Faverjon--Roques; no
support-growth, convergence or solvability theorem for arbitrary
variable-coefficient Mahler systems is inferred
(Section~\ref{dsup:sec:intro}, Section~\ref{dsup:sub:notmade}).
\item No claim that these structures satisfy the residue difference
closedness hypotheses of Pal's relative quantifier elimination, and the
coefficient-fixing structure is not identified with the existentially
closed residue difference structures of other theories of valued
difference fields (Section~\ref{dsup:sec:intro}, the remark after
Theorem~\ref{dsup:thm:tp2}).
\item The coefficient action is the identity; a nontrivial difference
automorphism of $\kk$ changes the fixed-exponent recurrences
(Section~\ref{dsup:sub:notmade}, Section~\ref{dsup:sec:questions}).
\item Finite commuting families only; no contraction is asserted for a
nonabelian acting group (Section~\ref{dsup:sub:notmade}).
\item The multiplicative results concern the abelian group $K^\times$, not
matrix gauge classification; no universal smaller formula for
$H^r(\Z^d,M)$ is claimed for arbitrary shears and characters
(Sections~\ref{dsup:sec:multiplicative} and~\ref{dsup:sub:notmade}).
\item The resonance projection is linear but not a field homomorphism, and
$[x]_0$ is not a ring homomorphism on $K$ (Section~\ref{dsup:sec:setup}, the
remark after Corollary~\ref{dsup:cor:poly}).
\item The kernel statements do not say that $K$ is a direct sum of
eigenspaces (after Corollary~\ref{dsup:cor:fixed}).
\item The Koszul complex is a complex of additive groups; its naive
coefficientwise wedge product is not claimed to make $\dd$ an ordinary
derivation, and $\Delta_j$ satisfies no ordinary Leibniz rule
(Section~\ref{dsup:sec:koszul}).
\item Lemma~\ref{dsup:lem:orbit} is not an assertion of cofinal growth;
strong summation is not valuation convergence of finite partial sums, and
no support proof is replaced by ``the orbit tends to infinity''
(Sections~\ref{dsup:sec:setup} and~\ref{dsup:sec:orbits}, and the shear
example~\eqref{dsup:eq:shear}).
\item The single-equation criteria \eqref{dsup:eq:scalaradd} and
\eqref{dsup:eq:scalarmult} are existence equivalences, not claims that the
equations are solvable for every right-hand side
(Section~\ref{dsup:sec:dilations}).
\item Full Hahn fields are essential: restricted support classes need
separate closure checks, and definability in the full field supplies no
algorithm for arbitrary presentations (the remark on fullness in
Section~\ref{dsup:sec:definability}, Section~\ref{dsup:sub:notmade},
Section~\ref{dsup:sec:questions}).
\item Characteristic zero is retained; the monomial-recognition formula is
not exported to characteristic two (Section~\ref{dsup:sub:notmade}).
\item The undecidability theorem concerns the complete theory of the
specified structure; it asserts no failure of decidability for every
smaller language of a valued field, and interpretation with a named
parameter alone would not justify the reduct claim
(Section~\ref{dsup:sec:logic}).
\item The interpreted arithmetic does not follow merely from naming a
valuation ring; adding $S_2$ is a genuine expansion, so there is no
contradiction with tame valued-field reducts (the remark after
Theorem~\ref{dsup:thm:tp2}).
\item Interdefinability of the rational dilations does not imply
conjugacy (Section~\ref{dsup:sec:alldilations}).
\item Ordinary group cohomology and complete first-order theories are
set-sized only. There is no set of class automorphisms or class cochains;
the centralizer of class automorphisms is structural shorthand; the
support localizations carry no birthday bound, and no birthday fragment is
asserted closed under them (Section~\ref{dsup:sec:surreal}).
\item On $\No(i)$ the real axis $\No$ is not recovered, and the centralizer
contains all automorphisms of $\C$; preserving the recovered valuation
structure is not preserving a chosen real form
(Section~\ref{dsup:sub:recovered}).
\item $S_q$ is not the surreal exponential, the omega-map, a
Berarducci--Mantova derivation or an automorphism known to preserve any of
these; that strong coefficient-fixing automorphisms commute with the
infinitesimal Hahn exponential is not a claim that a monomial automorphism
commutes with the global surreal exponential (Sections~\ref{dsup:sec:multiplicative} and~\ref{dsup:sub:recovered}).
\item The fine topology of the full surreal field is not the summation
mechanism of any proof (Section~\ref{dsup:sub:recovered}).
\item No uniform algorithm for arbitrary Hahn names is supplied; deciding a
movement chamber or testing coefficient equality may need unavailable
information (Section~\ref{dsup:sub:notmade}).
\item The questions of Section~\ref{dsup:sec:questions} are raised here, not
claimed to have been posed before in exactly this form, and the boundary of
the method they describe is not a claim that the general problem is
impossible (Section~\ref{dsup:sec:questions}).
\item The 66,131 finite checks are not a formal verification: they do not
decide Hahn summability, quantify over ordered groups, check the
first-order interpretations or verify proper-class assertions; compiling
and inspecting the PDF and running the script were not independent review
(title page, Appendices~\ref{dsup:app:verification}
and~\ref{dsup:app:provenance}; the delivered README and
\repo{PROOF_AUDIT.md}).
\item Added: none of the four questions of the surcomplex-automorphisms
report (\lbl{saut:sec:questions}) is answered. Its non-claims that the
four-layer decomposition covers $G_v$ only and that $\Aut(\No(i))$ is not
classified are extended only by the one subgroup $\Cen(S_q)$, which is
shown to lie in $G_v$; the unrestricted group is not classified
(Section~\ref{dsup:sec:position}).
\item Added: the rigidity report's Question~13.1 (\lbl{q:image}) is
untouched; nothing here says which canonical lifts $\widehat T$ commute with
the exponential (Section~\ref{dsup:sec:position}).
\item Added: the dilations here are not the argument dilations
$f(z)\mapsto f(\lambda z)$ of the holonomic-rigidity report, and no
statement of that report is used or extended
(Section~\ref{dsup:sec:position}).
\item Added: the arithmetic interpretation and undecidability of the
birthday-cutoffs report concern a different expansion of $\No$; neither
result implies the other (Section~\ref{dsup:sec:position}).
\item Added: no \texttt{dsup:} label has a Lean mapping in
\repo{docs/FORMALIZATION.md}; the Lean Hahn-series interfaces the source
cites (strong summation, the Neumann lemma, infinitesimal exponential and
logarithm, algebraic closedness) prove none of this report's statements
(Section~\ref{dsup:sec:position}, Appendix~\ref{dsup:app:pinned}).
\end{enumerate}

\section{Further questions suggested by the results}\label{dsup:sec:questions}

The following are questions raised by this article, not problems claimed
to have been previously posed in exactly this form.

\begin{question}[Minimal closure for definable normal forms]
Characterize $S_2$-invariant subfields $L\subseteq\kk((t^\Gamma))$ on which
$\im(S_2-I)=\{f\in L:[f]_0=0\}$ and every required coefficient shift
$x/t^g$ remains available. Which natural restricted-series fields recover
truncation without containing all full Hahn subseries?
\end{question}

The proof shows what must be tested: the two half-orbit sums supply
additive witnesses, while truncation additionally requires the relevant
subseries to exist in $L$. Those are different closure requirements.
Failure of one need not imply failure of the other.

\begin{question}[Nontrivial coefficient difference structure]
Let $\rho\in\Aut(\kk)$ and
$\sigma(\sum a_gt^g)=\sum\rho(a_g)\chi(g)t^{\theta g}$. Can one obtain an
equally explicit decomposition in which each fixed-exponent coefficient
is replaced by the kernel and cokernel of a scalar-twisted operator on
$(\kk,\rho)$?
\end{question}

The increasing-orbit support lemma still applies. The obstacle is not
support transport but the residual coefficient equations and their
compatibility under several automorphisms. A satisfactory result would
separate those two mechanisms rather than impose residue difference
closedness as an unexplained global substitute.

\begin{question}[Definability from other exponent automorphisms]
Which ordered exponent automorphisms $\theta$ admit a finite collection
of polynomial equations in $x,\sigma x,\ldots$ that recognizes canonical
monomials? Can the centralizer argument be extended to shears with a
large fixed Hahn field, where the definable projection lands in that
field rather than in $\kk$?
\end{question}

For rational dilations, the power equation and the vanishing of the
relevant scalar eigenspace solve both tasks. A shear lacks that simple
power relation; its fixed-field projection alone does not distinguish
ordinary coefficients from all fixed series. This is a concrete boundary
of the method, not a claim that the general problem is impossible.

\section{Conclusion}

The article begins with a support calculation and ends with a structural
and logical reconstruction theorem. The support calculation has no
rank-one hypothesis: moving exponent orbits have summable increasing
halves, and only pointwise-fixed resonant exponents obstruct scalar
difference equations. For commuting actions, the normalized inverses
combine into an explicit contraction in every cohomological degree.
The multiplicative problem then separates into a coupled exponent/
leading-coefficient condition and an exactly computable logarithmic
obstruction.

For doubling, the resonant projection is just the ordinary constant
coefficient. Together with the equation $S_2x=x^2$, it recovers the entire
normal form inside the field language with one automorphism. This forces
all centralizing field automorphisms to be strong untwisted lifts, and
it makes full arithmetic support coding available. In this precise
sense, a single elementary symmetry both rigidifies field automorphisms
and greatly increases the information definable in the field. The same
proofs yield explicit new statements for surreal and surcomplex normal
forms without confusing proper-class topology, formal Hahn summation,
or the separate exponential structure.

\appendix
\section{Endpoint identities and finite verification}\label{dsup:app:verification}

The accompanying \texttt{verify.py} uses only Python's standard library
and exact rational arithmetic. Its output is recorded in
\texttt{verification.txt}. (As placed in the collection these are
\repo{code/verify.py} and \repo{data/verification.txt};
Appendix~\ref{dsup:app:files} gives the rerun.) It is intentionally a collection of finite
checks rather than a finite approximation advertised as proof of an
infinite statement.

For a truncation of the increasing-chamber part, define
\[
 Q^+_{\lambda,N}f=-\sum_{n=0}^{N-1}
                \lambda^{-n-1}\sigma^nf.
\]
A finite telescope gives
\begin{equation}\label{dsup:eq:finiteplus}
 (\sigma-\lambda)Q^+_{\lambda,N}f
   =f-\lambda^{-N}\sigma^Nf.
\end{equation}
For the decreasing chamber, define
\[
 Q^-_{\lambda,N}f=\sum_{n=1}^{N}
                  \lambda^{n-1}\sigma^{-n}f.
\]
Then
\begin{equation}\label{dsup:eq:finiteminus}
 (\sigma-\lambda)Q^-_{\lambda,N}f
   =f-\lambda^N\sigma^{-N}f.
\end{equation}
The program tests these identities on rational exponents under doubling
and on integer-coordinate exponents under shears, including nontrivial
characters. It retains the endpoint term exactly rather than treating it
as a small numerical error.

The program also checks the Koszul identities on finite diagonal
coefficient models in every exterior degree for $d=1,\ldots,5$. These
models exercise the signs of $\eps_j,\iota_j$ and the ordered factors
$E_{j-1}$, including coordinate directions with zero difference operator.
It checks both the homotopy identity and the annihilation/normalization
identities on basis vectors. This is a check of the finite operator
algebra appearing in the proof, not a reduction of arbitrary Hahn
operators to finite diagonal matrices.

Additional checks compare the first finitely many coefficients in the
multiplicative product of Example~\ref{dsup:ex:product},
test explicit phase-twist resonance projections, and demonstrate the
failure of a nontrivial character twist to commute with doubling on a
monomial. The program does not decide Hahn summability, quantify over
all ordered groups, check the first-order interpretations, or verify
proper-class assertions. Those tasks remain in the written proofs.

\section{A compact formula dictionary}\label{dsup:app:formulas}

This dictionary uses $S_2$ and the definable functions whose existence
and uniqueness were proved above. Every occurrence of a function can
be eliminated in favor of its graph formula.

\renewcommand{\arraystretch}{1.25}
\begin{longtable}{@{}p{.22\textwidth}p{.73\textwidth}@{}}
\toprule
Recovered object & Definition or test \\
\midrule\endhead
Coefficient field & $C(c)\iff S_2c=c$.\\
Constant coefficient & $B(x)=c\iff C(c)\land\exists y\,(x-c=S_2y-y)$.\\
Canonical monomial & $M(m)\iff m\ne0\land S_2m=m^2$.\\
Coefficient at $m$ & $\coef(x,m)=B(x/m)$, for $M(m)$.\\
Support & $\Supp(m,x)\iff M(m)\land\coef(x,m)\ne0$.\\
Positive exponent & $J(m)\iff M(m)\land m\ne1\land B((1-m)^{-1})=1$.\\
Exponent comparison & $v(m)<v(n)\iff J(n/m)$, for monomials.\\
Valuation ring & Every monomial in the support is $1$ or satisfies $J$.\\
Infinitesimal ideal & Every monomial in the support satisfies $J$.\\
Leading monomial & The unique support monomial of least recovered exponent.\\
Truncation at $m$ & Preserve the coefficient at $n$ exactly when $J(m/n)$; set it to zero otherwise.\\
Natural-number sort & $N_m=\{u:\Supp(u,(1-m)^{-1})\}$, for $J(m)$.\\
Power-set sort & Series with coefficients in $\{0,1\}$ and support contained in $N_m$.\\
\bottomrule
\end{longtable}

The monomial $1$ has exponent zero, not natural-number index one. In the
interpreted natural-number sort the field element $1=m^0$ represents
$0\in\N$, whereas the parameter $m$ represents $1\in\N$.
This distinction is essential when translating arithmetic formulas.

\section{Provenance, reproducibility, and novelty ledger}\label{dsup:app:provenance}

The delivery contained this standalone \LaTeX{} source, its compiled PDF,
a standard-library verification script, the exact output of that script,
a README, and a small Makefile. The bibliography is internal, so no
Bib\TeX{} database is required. In the delivered flat layout a typical
build was
\begin{lstlisting}
python3 verify.py
latexmk -pdf -interaction=nonstopmode -halt-on-error article.tex
\end{lstlisting}
The PDF was compiled from the delivered source and visually inspected
through rendered page images. The finite test output was generated by
running the delivered script. Neither activity is described as an
independent mathematical review. (As placed in the collection the delivered
PDF and README are not shipped; the files and working commands are in
Appendix~\ref{dsup:app:files}.)

\begin{longtable}{@{}>{\raggedright\arraybackslash}p{.27\textwidth}>{\raggedright\arraybackslash}p{.68\textwidth}@{}}
\toprule
Result or ingredient & Status in this manuscript \\
\midrule\endhead
Hahn field operations; normal forms & Standard background, with the normal-form dictionary cited to the literature.\\
Monomial and character lifts & Prior construction; explicitly credited to the literature and repository.\\
Orbit-resolvent theorem & Proved here for arbitrary ordered exponent automorphisms; proposed addition, priority not certified.\\
Full Koszul contraction & Proved here with an explicit strong homotopy; proposed addition.\\
Principal-unit logarithm & Standard formal Hahn construction; used, not claimed new.\\
Multiplicative obstruction theorem & Proved here from the contraction and multiplicative coordinates; proposed addition in this general form.\\
Single-dilation recovery & Main proposed new structural result; explicit formulas and proofs supplied.\\
Unrestricted centralizer & Main proposed new algebraic consequence; no valuation or strength hypothesis assumed.\\
Support-to-arithmetic principle & Prior precedent in Camacho; a Boolean coding and uniform parameter elimination are provided in the present expansion.\\
Rational interdefinability and nonconjugacy & Proved consequences with distinct meanings; not a claim about exponential automorphisms.\\
Lean coverage & None claimed for these new theorems. A dependency route, not a completed formalization, is supplied.\\
Published open problems & No named existing open problem is claimed settled; limitations are explicit.\\
\bottomrule
\end{longtable}

The bibliography review was focused on Hahn automorphisms, Mahler
solutions, valued difference fields, and truncation definability. It was
not an exhaustive search of all unpublished notes or every language's
literature. A subsequently located antecedent would change the priority
assessment, not the logical content of the proofs. The intended research
contribution is the collection of exact statements and their deductions,
with the single-dilation recovery theorem as its organizing result.

\section{Provenance of this report}\label{dsup:app:report}

\subsection{One manuscript}\label{dsup:app:onesource}

This report is built from \emph{one} manuscript: the archive
\texttt{surreal\_dilation\_support}, manuscript~09 of batch~22 of the
collection, dated 22 September 2026 and pinned to repository commit
\nolinkurl{465a54b479a1ee842cbf7db1689a7d2f6bfe25e1}. It was placed in commit
\texttt{7b5f934} as a new report: it answers no question that a report of
the collection names, and it continues only one non-claim of the
surcomplex-automorphisms report (Section~\ref{dsup:sec:position}). There was
no second source; nothing was merged and nothing was selected away. Every
result, proof, example, question and limitation of the manuscript is
printed here.

The delivered verification program \repo{code/verify.py}, its Makefile
\repo{code/Makefile}, the recorded output \repo{data/verification.txt} and
the author-side audit \repo{PROOF_AUDIT.md} are kept byte-identical. The
audit and the Makefile describe the delivered flat layout (see
Appendix~\ref{dsup:app:files}). The manuscript's article and README were
rewritten, and the delivered PDF was replaced by a build of this text.

The editorial changes are these.
\begin{enumerate}[label=(\arabic*)]
\item Every label carries the prefix \texttt{dsup:}. The source's 99 labels
are kept, unchanged after the prefix, and nine were added
(\texttt{dsup:sec:conventions}, \texttt{dsup:sec:position},
\texttt{dsup:sub:recovered}, \texttt{dsup:sub:notmade},
\texttt{dsup:sec:nonclaims}, \texttt{dsup:app:report},
\texttt{dsup:app:onesource}, \texttt{dsup:app:pinned},
\texttt{dsup:app:files}).
\item The exponent dilation $D_q$ is written $S_q$, the abbreviation $D$ for
$D_2$ is dropped, and the support formula $S(m,x)$ is written $\Supp(m,x)$
(Section~\ref{dsup:sec:conventions}). The sentence introducing the language
at the start of Section~\ref{dsup:sec:definability} and the first sentence of
Appendix~\ref{dsup:app:formulas} were adjusted accordingly.
\item The title page names the report as a single-source report of the
collection and extends the status paragraph. Sections~\ref{dsup:sec:conventions}
and~\ref{dsup:sec:position}, the ledger of Section~\ref{dsup:sec:nonclaims}
(with five items added) and this appendix are new. Short sentences pointing
to them, and to the files as placed, were added at the end of
Section~\ref{dsup:sec:intro} and in Appendices~\ref{dsup:app:verification}
and~\ref{dsup:app:provenance}; the latter's description of the delivery is
put in the past tense.
\item The columns of the two appendix tables are set ragged right, which
removes two underfull lines of the delivered build.
\end{enumerate}
No mathematical statement was changed. Where the edit had to choose: the
dilation was renamed rather than keeping $D_q$, because the notation guide
reserves $D$ for a scalar derivation and the nearest report already calls
the map $S_a$; the support formula was renamed rather than giving the
dilation a third letter; the title is kept; and the report is placed among
the surcomplex reports, because its results are uniform in the coefficient
field, Theorem~\ref{dsup:thm:surreal} treats $\No$ and $\No(i)$ together,
and its nearest neighbour is the surcomplex-automorphisms report.

\subsection{Corrections since the pin}\label{dsup:app:pinned}

The source compared itself with the collection at its pin. Checked against
the collection at the placement commit \texttt{7b5f934}, its repository
statements stand; one addition matters.
\begin{itemize}[leftmargin=*,itemsep=3pt]
\item The root README still documents strong Hahn summation, the full
Neumann support lemma, coefficient maps preserving strong sums, the
infinitesimal exponential and logarithm, and algebraic closedness of Hahn
fields over algebraically closed characteristic-zero coefficients; these
are Lean statements of the collection's Hahn layer, and none of them is a
statement of this report. The catalogue still separates the small-net
behaviour of the fine topology from strong Hahn summation. The
surcomplex-automorphisms report still constructs the coefficient, exponent
and character lifts (\lbl{saut:thm:monomial}); neither its article nor that
of the rigidity report has changed since the pin.
\item At the pin the catalogue listed 44 research reports in five families.
Batch~21, placed in \texttt{d4e71b7} after the pin, added two, among them
the birthday-cutoffs report whose undecidability results are compared in
Section~\ref{dsup:sec:position}; batch~22, placed in \texttt{7b5f934}, adds
further reports, this one among them.
\item The surcomplex-automorphisms guide that the source inspected gained,
after the pin, one paragraph on the two batch-21 reports that use
conjugation on $\No[i]$; it bears on no statement here. The source did not
audit that report's article; the relations of
Section~\ref{dsup:sec:position} were checked against the article itself.
\end{itemize}

\subsection{Files as placed and rerun}\label{dsup:app:files}

\begin{center}\small
\begin{tabular}{@{}ll@{}}
\toprule
File & Role \\
\midrule
\repo{code/verify.py} & the delivered finite checks (delivered as \texttt{verify.py})\\
\repo{code/Makefile} & the delivered Makefile, for the flat delivered layout\\
\repo{data/verification.txt} & the delivered recorded output\\
\repo{PROOF_AUDIT.md} & the delivered author-side audit\\
\bottomrule
\end{tabular}
\end{center}

The program uses only the Python standard library, prints to standard
output and writes no file. The Makefile, however, assumes that
\texttt{article.tex}, \texttt{verify.py} and \texttt{verification.txt} sit in
one directory: its target \texttt{check} runs
\texttt{python3 verify.py > verification.txt}, overwriting the recorded output
in place, and \texttt{all} also rebuilds \texttt{article.pdf}. It must not
be run in this directory. Two working reruns, from the report directory,
with a fresh scratch directory \texttt{\$T}:
\begin{lstlisting}
python code/verify.py > $T/direct.txt
diff --strip-trailing-cr $T/direct.txt data/verification.txt

cp code/verify.py code/Makefile $T/ && make -C $T check
diff --strip-trailing-cr $T/verification.txt data/verification.txt
\end{lstlisting}
Both were run for this report under Python~3.14.4 on Windows (the second
with GNU Make). Each printed 17 \texttt{PASS} groups and 66,131 checks and
reproduced \repo{data/verification.txt} exactly apart from line endings
(Windows writes CRLF).

\begin{thebibliography}{99}
\small
\raggedright

\bibitem{KKS}
Elliot Kaplan, Lothar Sebastian Krapp, and Michele Serra,
\emph{Decomposing the automorphism group of the surreal numbers},
arXiv:2509.22374, version 3, 23 April 2026.
\url{https://arxiv.org/abs/2509.22374v3}.
The normal-form and monomial-automorphism preliminaries, and the distinction
from exponential-preserving automorphisms, are the relevant background.

\bibitem{Camacho}
Santiago Camacho,
\emph{Truncation in Hahn Fields is Undecidable and Wild},
arXiv:1706.03722, version 1, 12 June 2017.
\url{https://arxiv.org/abs/1706.03722v1}.
Sections 4--6 provide the arithmetic-interpretation and support-definability
precedent. The present article adds a distinguished dilation instead of
assuming the same truncation signature.

\bibitem{FR}
Colin Faverjon and Julien Roques,
\emph{Hahn series and Mahler equations: algorithmic aspects},
Journal of the London Mathematical Society \textbf{110} (2024), e12945.
\url{https://doi.org/10.1112/jlms.12945}.
Accepted manuscript: arXiv:2412.04928v1,
\url{https://arxiv.org/abs/2412.04928v1}.

\bibitem{Pal}
Koushik Pal,
\emph{Multiplicative Valued Difference Fields},
arXiv:1011.1655, version 2, 13 November 2010.
\url{https://arxiv.org/abs/1011.1655v2}.

\bibitem{Repo}
Vladimir Reshetnikov, \emph{Surreal} repository, root README and documented
Hahn-series formalization interfaces, snapshot
\texttt{465a54b479a1ee842cbf7db1689a7d2f6bfe25e1}, reviewed
22 September 2026.
\url{https://github.com/VladimirReshetnikov/Surreal/tree/465a54b479a1ee842cbf7db1689a7d2f6bfe25e1}.

\bibitem{RepoDocs}
Vladimir Reshetnikov, \emph{The surreal and surcomplex reports},
\texttt{docs/README.md} in the same repository snapshot.
AI-assisted research-report catalogue and statement of verification scope.
\url{https://github.com/VladimirReshetnikov/Surreal/blob/465a54b479a1ee842cbf7db1689a7d2f6bfe25e1/docs/README.md}.

\bibitem{RepoAut}
\emph{Automorphisms of the Surcomplex Numbers: Real forms, Hahn symmetries,
valuation, and rigidity}, research report prepared with ChatGPT,
22 September 2026; repository guide
\texttt{docs/surcomplex/surcomplex-field-automorphisms/README.md},
in the same snapshot. The guide's statement inventory was inspected;
its complete article was not independently audited here.
\url{https://github.com/VladimirReshetnikov/Surreal/tree/465a54b479a1ee842cbf7db1689a7d2f6bfe25e1/docs/surcomplex/surcomplex-field-automorphisms}.

\end{thebibliography}
\end{document}
